\chapter[Fans and Toric Varieties]
{Fans and Toric Varieties:\\ Gluing Affine Pieces}
\label{chap:fans-and-toric-varieties}

\begin{chapteroverview}
A single strongly convex rational cone
\[
\sigma\subseteq N_{\RR}
\]
produces one normal affine toric variety
\[
U_\sigma
=
\Spec\Bbbk[\sigma^\vee\cap M].
\]
But many of the toric varieties we care about most—projective spaces,
Hirzebruch surfaces, and blowups—are not affine. They cannot be
described by one cone.

The solution is local-to-global. Several cones will provide several
affine toric charts, and common faces will provide the affine open
subsets along which those charts are glued. A compatible collection
of cones is called a \emph{fan}.

This chapter develops the passage
\[
\begin{array}{c}
\text{cones and common faces}
\\[-1mm]
\downarrow
\\[-1mm]
\text{affine toric charts and common open subsets}
\\[-1mm]
\downarrow
\\[-1mm]
\text{a global toric variety }X_\Sigma.
\end{array}
\]
The running examples will be toric curves and surfaces, especially
\(\PP^1\), \(\PP^2\), products, and blowups.
\end{chapteroverview}

\begin{threedictionaries}
The global theory adds one new operation to each dictionary.

\medskip

\begin{center}
\small
\begin{tabularx}{\linewidth}{@{}
  >{\raggedright\arraybackslash}X
  >{\raggedright\arraybackslash}X
  >{\raggedright\arraybackslash}X@{}}
\toprule
\textbf{Algebra}
&
\textbf{Geometry}
&
\textbf{Combinatorics}
\\
\midrule
semigroup algebras
\(\Bbbk[\sigma^\vee\cap M]\)
&
affine charts \(U_\sigma\)
&
cones \(\sigma\subseteq N_{\RR}\)
\\[1mm]
localization at characters
&
common affine open subsets
&
common faces
\\[1mm]
compatible restriction maps
&
gluing affine charts
&
a fan \(\Sigma\)
\\
\bottomrule
\end{tabularx}
\end{center}
\end{threedictionaries}

\section{Why one cone is not enough}
\label{sec:why-one-cone-not-enough}
\index{affine chart!toric}
\index{gluing!affine toric charts}

Every cone gives an affine variety. This is a strength of the
construction
\[
\sigma\longmapsto
U_\sigma
=
\Spec\Bbbk[\sigma^\vee\cap M],
\]
but it is also a limitation. Neither \(\PP^1\) nor \(\PP^2\) is
affine, so neither can come from a single cone.

They are, however, covered by affine toric varieties. The question is
therefore not how to replace cones, but how to use several of them at
once.

\begin{guidingquestion}
How can several cones encode affine toric charts together with the
instructions for gluing those charts?
\end{guidingquestion}

\subsection{The affine construction is a local model}

Let
\[
\tau\preceq\sigma
\]
be a face. Chapter~2 showed that
\[
U_\tau\subseteq U_\sigma
\]
is a torus-invariant affine open subset. More precisely, if
\[
\tau=\sigma\cap m^\perp
\qquad
\text{for }m\in\sigma^\vee\cap M,
\]
then
\[
U_\tau
=
D(\character{m})
\subseteq U_\sigma.
\]

Thus a face already carries gluing information:

\[
\boxed{
\text{a face of a cone}
\quad\longleftrightarrow\quad
\text{an affine open subset of the corresponding chart}.
}
\]

Suppose now that two cones \(\sigma_1,\sigma_2\subseteq N_{\RR}\)
meet in a cone
\[
\tau=\sigma_1\cap\sigma_2
\]
that is a face of both. Then the same affine toric variety \(U_\tau\)
appears as an open subset of both:
\[
U_{\sigma_1}
\supseteq
U_\tau
\subseteq
U_{\sigma_2}.
\]
The identity on the common dense torus identifies these two copies of
\(U_\tau\). This gives exactly the kind of overlap needed to glue the
two affine charts.

We first watch this happen in dimension one.

\subsection{\texorpdfstring{\(\PP^1\)}{P1} from two affine lines}

Take
\[
N=\ZZ,
\qquad
M=\ZZ,
\]
and let \(t=\character{1}\) be the coordinate on the torus
\[
T=\Gm.
\]
Consider the two rays
\[
\sigma_+=\RR_{\geq0},
\qquad
\sigma_-=\RR_{\leq0},
\]
whose intersection is the zero cone
\[
\tau=\{0\}.
\]

Their dual semigroups are
\[
\sigma_+^\vee\cap M=\NN,
\qquad
\sigma_-^\vee\cap M=-\NN,
\qquad
\tau^\vee\cap M=\ZZ.
\]
Consequently,
\[
U_{\sigma_+}
=
\Spec\Bbbk[t]
\simeq\AAA^1,
\]
\[
U_{\sigma_-}
=
\Spec\Bbbk[t^{-1}]
\simeq\AAA^1,
\]
and
\[
U_\tau
=
\Spec\Bbbk[t,t^{-1}]
=
\Gm.
\]

The face inclusions identify \(U_\tau\) with
\[
D(t)\subseteq U_{\sigma_+}
\]
and
\[
D(t^{-1})\subseteq U_{\sigma_-}.
\]
If we write
\[
x=t
\qquad\text{on }U_{\sigma_+},
\]
and
\[
y=t^{-1}
\qquad\text{on }U_{\sigma_-},
\]
then the transition function on the overlap is
\[
y=x^{-1}.
\]

\begin{figure}[ht]
\centering
\begin{tikzpicture}[
  chart/.style={
    rounded corners,
    draw=black!55,
    minimum width=3.1cm,
    minimum height=1.35cm,
    align=center,
    inner sep=5pt
  }
]
  % fan
  \begin{scope}[xshift=-4.8cm]
    \draw[very thick,DictionaryBlue,-{Latex[length=2.3mm]}]
      (0,0)--(2.4,0);
    \draw[very thick,SurfaceGreen,-{Latex[length=2.3mm]}]
      (0,0)--(-2.4,0);
    \fill[WarningRed] (0,0) circle (2.7pt);
    \node[DictionaryBlue,above] at (1.45,0)
      {\(\sigma_+\)};
    \node[SurfaceGreen,above] at (-1.45,0)
      {\(\sigma_-\)};
    \node[WarningRed,below] at (0,-0.05)
      {\(\tau=\{0\}\)};
    \node[font=\bfseries] at (0,1.15)
      {cones in \(N_{\RR}\)};
  \end{scope}

  % charts
  \node[chart,fill=DictionaryBlue!7] (plus) at (2.2,1.1)
    {\(U_{\sigma_+}\simeq\AAA^1_x\)};
  \node[chart,fill=SurfaceGreen!7] (minus) at (2.2,-1.1)
    {\(U_{\sigma_-}\simeq\AAA^1_y\)};
  \node[
    rounded corners,
    draw=PredictionTeal,
    fill=PredictionTeal!5,
    align=center,
    minimum width=2.5cm,
    minimum height=1.0cm
  ] (torus) at (6.0,0)
    {\(U_\tau=\Gm\)\\ \(y=x^{-1}\)};

  \draw[very thick,-{Latex[length=2.3mm]},DictionaryBlue]
    (plus.east)--(torus.north west);
  \draw[very thick,-{Latex[length=2.3mm]},SurfaceGreen]
    (minus.east)--(torus.south west);
\end{tikzpicture}
\caption{The two rays give affine-line charts, while their common
zero face gives the torus on which the charts overlap.}
\label{fig:p1-two-cones-two-charts}
\end{figure}

Gluing these two affine lines along \(\Gm\) using
\[
y=x^{-1}
\]
produces \(\PP^1\). Indeed, the maps
\[
\AAA^1_x\longrightarrow\PP^1,
\qquad
x\longmapsto[1:x],
\]
and
\[
\AAA^1_y\longrightarrow\PP^1,
\qquad
y\longmapsto[y:1],
\]
agree on the overlap precisely because \(y=x^{-1}\).

\noindent
\textit{Reference.}
This is the one-dimensional case of Cox--Little--Schenck,
Examples~3.1.10 and~3.1.11. The general affine gluing model for
projective space is developed in \S3.0.

\begin{pauseandthink}
The point \(0\in\AAA^1_x\) becomes \([1:0]\), while
\(0\in\AAA^1_y\) becomes \([0:1]\). Why are these two points not part
of the overlap? Explain how the two origins supply the two boundary
points of
\[
\PP^1\setminus\Gm.
\]
\end{pauseandthink}

This first example contains the global construction in miniature:

\[
\begin{array}{c@{\qquad\longleftrightarrow\qquad}c}
\sigma_+,\sigma_-
&
\AAA^1_x,\AAA^1_y
\\[1mm]
\sigma_+\cap\sigma_-=\{0\}
&
\AAA^1_x\cap\AAA^1_y=\Gm.
\end{array}
\]

\subsection{Common faces are the transition charts}

The gluing in the previous example was not chosen independently of
the cones. It was forced by the shared face.

For a face
\[
\tau\preceq\sigma,
\]
the inclusion
\[
U_\tau\hookrightarrow U_\sigma
\]
restricts to the identity on the dense torus \(T\). Hence whenever
\[
\tau=\sigma_1\cap\sigma_2
\]
is a face of both cones, the diagram
\[
\begin{tikzcd}[column sep=large]
&
U_\tau
\arrow[dl,hook']
\arrow[dr,hook]
&
\\
U_{\sigma_1}
&&
U_{\sigma_2}
\end{tikzcd}
\]
provides canonical gluing data.

On coordinate rings, these open embeddings come from localizations.
If
\[
\tau
=
\sigma_1\cap m^\perp
=
\sigma_2\cap(-m)^\perp
\]
for a suitable \(m\in M\), then
\[
\Bbbk[U_\tau]
=
\Bbbk[U_{\sigma_1}]_{\character{m}}
=
\Bbbk[U_{\sigma_2}]_{\character{-m}}.
\]
The transition map is the identity inside the common Laurent
polynomial ring \(\Bbbk[M]\).

\begin{prediction}
Suppose a finite collection of cones is meant to describe affine
charts of one global toric variety.

What conditions should be imposed so that every boundary piece needed
for an overlap is present, and so that the intersection of any two
charts is represented by one common face rather than by incompatible
pieces?
\end{prediction}

\subsection{\texorpdfstring{\(\PP^2\)}{P2} from three affine planes}

Let
\[
N=\ZZ^2,
\qquad
e_0=-e_1-e_2,
\]
and consider
\[
\sigma_0=\Cone{e_1,e_2},
\qquad
\sigma_1=\Cone{e_0,e_2},
\qquad
\sigma_2=\Cone{e_1,e_0}.
\]

\begin{figure}[ht]
\centering
\begin{tikzpicture}[scale=1.1]
  \fill[DictionaryBlue!9]
    (0,0)--(2.35,0)--(0,2.35)--cycle;
  \fill[SurfaceGreen!9]
    (0,0)--(0,2.35)--(-1.8,-1.8)--cycle;
  \fill[ExampleGold!13]
    (0,0)--(-1.8,-1.8)--(2.35,0)--cycle;

  \draw[very thick,DictionaryBlue,-{Latex[length=2.4mm]}]
    (0,0)--(2.7,0);
  \draw[very thick,SurfaceGreen,-{Latex[length=2.4mm]}]
    (0,0)--(0,2.7);
  \draw[very thick,ExampleGold,-{Latex[length=2.4mm]}]
    (0,0)--(-2.15,-2.15);

  \node[DictionaryBlue,below] at (2.55,0) {\(e_1\)};
  \node[SurfaceGreen,left] at (0,2.55) {\(e_2\)};
  \node[ExampleGold,above left] at (-2.0,-2.0)
    {\(e_0=-e_1-e_2\)};

  \node at (0.75,0.8) {\(\sigma_0\)};
  \node at (-0.75,0.65) {\(\sigma_1\)};
  \node at (0.2,-0.85) {\(\sigma_2\)};
\end{tikzpicture}
\caption{The three maximal cones for \(\PP^2\). Their rays and the
zero cone encode all chart intersections.}
\label{fig:p2-three-cones}
\end{figure}

Let \(m_1,m_2\) be dual to \(e_1,e_2\), and write
\[
x=\character{m_1},
\qquad
y=\character{m_2}.
\]
The complete affine atlas is contained in the following table.

\begin{center}
\footnotesize
\renewcommand{\arraystretch}{1.18}
\begin{tabularx}{\linewidth}{@{}
  >{\raggedright\arraybackslash}p{0.21\linewidth}
  >{\raggedright\arraybackslash}p{0.31\linewidth}
  >{\raggedright\arraybackslash}X@{}}
\toprule
\textbf{Cone or common face}
&
\textbf{Affine toric piece}
&
\textbf{Projective interpretation}
\\
\midrule
\(\sigma_0=\Cone{e_1,e_2}\)
&
\(\Bbbk[x,y]\), hence \(U_{\sigma_0}\simeq\AAA^2\)
&
\(V_0=\{z_0\neq0\}\);
\(\ (x,y)=(z_1/z_0,z_2/z_0)\)
\\
\(\sigma_1=\Cone{e_0,e_2}\)
&
\(\Bbbk[x^{-1},x^{-1}y]\), hence
\(U_{\sigma_1}\simeq\AAA^2\)
&
\(V_1=\{z_1\neq0\}\);
\(\ (x^{-1},x^{-1}y)=(z_0/z_1,z_2/z_1)\)
\\
\(\sigma_2=\Cone{e_1,e_0}\)
&
\(\Bbbk[y^{-1},xy^{-1}]\), hence
\(U_{\sigma_2}\simeq\AAA^2\)
&
\(V_2=\{z_2\neq0\}\);
\(\ (y^{-1},xy^{-1})=(z_0/z_2,z_1/z_2)\)
\\
\midrule
\(\sigma_0\cap\sigma_1=\Cone{e_2}\)
&
\(\Bbbk[x^{\pm1},y]\simeq
\Bbbk[x^{-1},x^{-1}y]_{x^{-1}}\)
&
\(V_0\cap V_1\);
\(\ (x,y)\mapsto(x^{-1},y/x)\)
\\
\(\sigma_0\cap\sigma_2=\Cone{e_1}\)
&
\(\Bbbk[x,y^{\pm1}]\)
&
\(V_0\cap V_2\);
\(\ (x,y)\mapsto(y^{-1},x/y)\)
\\
\(\sigma_1\cap\sigma_2=\Cone{e_0}\)
&
\(\Bbbk[a,b^{\pm1}]\), with
\(a=x^{-1}\), \(b=x^{-1}y\)
&
\(V_1\cap V_2\);
\(\ (a,b)\mapsto(a/b,1/b)\)
\\
\(\sigma_0\cap\sigma_1\cap\sigma_2=\{0\}\)
&
\(\Bbbk[x^{\pm1},y^{\pm1}]\), so \(U_{\{0\}}=\Gm^2\)
&
\(V_0\cap V_1\cap V_2\), the dense torus
\\
\bottomrule
\end{tabularx}
\end{center}

\noindent
\textit{Reference.}
This is Cox--Little--Schenck, Example~3.1.9. Its
\(n\)-dimensional generalization is Example~3.1.10.

\begin{surfacelaboratory}
Read the atlas directly from the table.

\begin{enumerate}[label=(\alph*),wide=0pt]
\item Verify that each maximal-cone ring is a polynomial ring in two
variables.

\item For each pair of charts, identify which character is inverted
and derive the displayed transition map from the two coordinate
systems.

\item Explain why every pairwise overlap is
\(\AAA^1\times\Gm\), while the triple overlap is \(\Gm^2\).

\item What two compatibility rules on a collection of cones are
suggested by this atlas? Your rules should ensure that all needed
overlap charts are present and that the intersection of any two
maximal charts is represented by one common cone.
\end{enumerate}
\end{surfacelaboratory}

The examples of \(\PP^1\) and \(\PP^2\) therefore lead directly to
the next definition: a compatible collection of cones should contain
all faces and should have the property that the intersection of any
two cones is a face of each.

\begin{guidingquestion}
How should we define such a collection precisely, and which features
of a global toric variety can already be read from the arrangement of
its cones?
\end{guidingquestion}


\section{Fans and their basic language}
\label{sec:fans-basic-language}
\index{fan!definition}
\index{fan!support}
\index{complete fan}
\index{subfan}

The examples of \(\PP^1\) and \(\PP^2\) forced two compatibility
conditions: faces must be retained, and intersections must be common
faces. We now make those conditions precise.

\subsection{The definition}

\begin{definition}[Fan]
\label{def:fan}
A \emph{fan} \(\Sigma\) in \(N_{\RR}\) is a finite collection of
cones such that:
\begin{enumerate}[label=(\alph*),wide=0pt]
\item every \(\sigma\in\Sigma\) is a strongly convex rational
polyhedral cone;
\item if \(\sigma\in\Sigma\) and \(\tau\preceq\sigma\), then
\(\tau\in\Sigma\);
\item if \(\sigma_1,\sigma_2\in\Sigma\), then
\[
\sigma_1\cap\sigma_2
\]
is a face of both \(\sigma_1\) and \(\sigma_2\).
\end{enumerate}
\end{definition}

\noindent
\textit{Reference.}
This is Cox--Little--Schenck, Definition~3.1.2.

The zero cone belongs to every nonempty fan, because it is a face of
every strongly convex cone.

The first condition ensures that every cone produces a normal affine
toric variety. The second ensures that every torus-invariant affine
open subset needed for gluing is represented. The third ensures that
two affine charts have one canonical common chart:
\[
\sigma_1\cap\sigma_2=\tau
\quad\leadsto\quad
U_{\sigma_1}\supseteq U_\tau\subseteq U_{\sigma_2}.
\]

\begin{example}[The affine fan of a cone]
Let \(\sigma\subseteq N_{\RR}\) be a strongly convex rational
polyhedral cone. Then
\[
\Sigma_\sigma
=
\{\tau:\tau\preceq\sigma\}
\]
is a fan. It consists of \(\sigma\) together with all of its faces.

The toric variety associated to this fan will be the affine variety
\(U_\sigma\).
\end{example}

\subsection{What can go wrong?}

Each part of \cref{def:fan} is necessary.

\begin{figure}[ht]
\centering
\begin{tikzpicture}[scale=0.84]
  % Good fan
  \begin{scope}[xshift=-6.0cm]
    \fill[DictionaryBlue!9]
      (0,0)--(2.2,0)--(0,2.2)--cycle;
    \fill[SurfaceGreen!9]
      (0,0)--(0,2.2)--(-1.7,-1.7)--cycle;
    \fill[ExampleGold!12]
      (0,0)--(-1.7,-1.7)--(2.2,0)--cycle;
    \draw[very thick,DictionaryBlue] (0,0)--(2.45,0);
    \draw[very thick,SurfaceGreen] (0,0)--(0,2.45);
    \draw[very thick,ExampleGold] (0,0)--(-1.95,-1.95);
    \fill (0,0) circle (2pt);
    \node[font=\bfseries] at (0,3.0) {a fan};
    \node[align=center] at (0,-2.65)
      {cones meet\\along common faces};
  \end{scope}

  % Missing faces
  \begin{scope}
    \fill[PastelYellow]
      (0,0)--(2.2,0)--(0,2.2)--cycle;
    \draw[very thick,WarningRed,dashed] (0,0)--(2.45,0);
    \draw[very thick,WarningRed,dashed] (0,0)--(0,2.45);
    \fill (0,0) circle (2pt);
    \node[font=\bfseries] at (0,3.0) {not a fan};
    \node[align=center] at (0,-2.65)
      {the cone is listed,\\but its rays are missing};
  \end{scope}

  % Bad intersection
  \begin{scope}[xshift=6.0cm]
    \fill[DictionaryBlue!11]
      (0,0)--(2.4,0)--(0,2.4)--cycle;
    \fill[SurfaceGreen!12,opacity=0.85]
      (0,0)--(1.7,1.7)--(-1.7,1.7)--cycle;
    \draw[very thick,DictionaryBlue] (0,0)--(2.55,0);
    \draw[very thick,DictionaryBlue] (0,0)--(0,2.55);
    \draw[very thick,SurfaceGreen] (0,0)--(1.9,1.9);
    \draw[very thick,SurfaceGreen] (0,0)--(-1.9,1.9);
    \draw[very thick,WarningRed] (0,0)--(1.65,1.65);
    \node[font=\bfseries] at (0,3.0) {not a fan};
    \node[align=center] at (0,-2.65)
      {the intersection is not\\a face of the blue cone};
  \end{scope}
\end{tikzpicture}
\caption{The face and intersection conditions distinguish fans from
arbitrary collections of cones.}
\label{fig:fan-and-two-nonfans}
\end{figure}

A cone containing a line also fails the definition: it is not strongly
convex, and the associated affine construction does not contain the
chosen torus as its dense torus in the required way.

\begin{pauseandthink}
Let
\[
\sigma_1=\Cone{(1,0),(0,1)}
\]
and
\[
\sigma_2=\Cone{(1,1),(-1,1)}.
\]
Compute \(\sigma_1\cap\sigma_2\). Why is it not a face of
\(\sigma_1\)? Which new ray would have to be inserted to subdivide
\(\sigma_1\) so that the two regions could become part of one fan?
\end{pauseandthink}

\subsection{Support and cones of each dimension}

\begin{definition}
Let \(\Sigma\) be a fan in \(N_{\RR}\).

Its \emph{support} is
\[
|\Sigma|
=
\bigcup_{\sigma\in\Sigma}\sigma
\subseteq N_{\RR}.
\]

For \(r\geq0\), we write
\[
\Sigma(r)
=
\{\sigma\in\Sigma:\dim\sigma=r\}.
\]
\end{definition}

\noindent
\textit{Reference.}
See Cox--Little--Schenck, immediately after Definition~3.1.2.

The set
\[
\Sigma(1)
\]
is the set of rays of the fan. For each ray \(\rho\in\Sigma(1)\), we
write
\[
u_\rho\in N
\]
for its primitive lattice generator.

A cone \(\sigma\in\Sigma\) is \emph{maximal} if it is not properly
contained in another cone of \(\Sigma\). We denote the set of maximal
cones by
\[
\Sigma_{\max}.
\]

Because the fan is finite, every cone lies in a maximal cone. Since
all faces of every maximal cone must be present, a fan is completely
determined by its maximal cones.

\begin{proposition}[Relative interiors partition the support]
\label{prop:relative-interiors-partition-fan}
For every fan \(\Sigma\),
\[
|\Sigma|
=
\bigsqcup_{\sigma\in\Sigma}
\relint(\sigma).
\]
Equivalently, every point of the support lies in the relative interior
of a unique cone of the fan.
\end{proposition}

\begin{proof}
Choose \(u\in|\Sigma|\), and let \(\sigma\in\Sigma\) contain \(u\).
Every point of a polyhedral cone lies in the relative interior of a
unique face of that cone. Hence
\[
u\in\relint(\tau)
\]
for a unique face \(\tau\preceq\sigma\), and \(\tau\in\Sigma\).

Suppose also that
\[
u\in\relint(\tau')
\]
for another cone \(\tau'\in\Sigma\). The intersection
\[
\tau\cap\tau'
\]
is a face of both. Since a relative-interior point cannot lie in a
proper face, we must have
\[
\tau=\tau\cap\tau'=\tau'.
\]
\end{proof}

Thus a fan is not merely a union of cones. It decomposes its support
into nonoverlapping pieces:
\[
\relint(\sigma),
\qquad
\sigma\in\Sigma.
\]

\begin{observation}
This partition is the combinatorial shadow of the orbit decomposition
that will appear later:
\[
\text{one cone}
\quad\longleftrightarrow\quad
\text{one torus orbit}.
\]
For now, it simply tells us that a direction \(u\in|\Sigma|\) belongs
unambiguously to one cone at the level of relative interiors.
\end{observation}

\subsection{The dimension of a fan}

We define
\[
\dim\Sigma
=
\max_{\sigma\in\Sigma}\dim\sigma
=
\dim|\Sigma|.
\]
The second equality holds because \(\Sigma\) is finite.

A fan is \emph{pure of dimension \(r\)} if every maximal cone has
dimension \(r\). The fans of \(\PP^1\) and \(\PP^2\) are pure of
dimensions \(1\) and \(2\), respectively.

\begin{warningbox}
The dimension of the fan need not equal the dimension of its toric
variety.

The fan
\[
\Sigma_{\Cone{(1,0)}}\subseteq\RR^2
\]
has dimension \(1\), but its affine toric variety is
\[
\AAA^1\times\Gm,
\]
which has dimension \(2\). The variety always has dimension
\(\rank N\); a lower-dimensional fan signals a torus factor.
\end{warningbox}

\subsection{Subfans}

\begin{definition}
A \emph{subfan} of \(\Sigma\) is a subset
\[
\Sigma'\subseteq\Sigma
\]
that is itself a fan.
\end{definition}

The phrase “itself a fan” matters. One may not remove a face while
retaining a cone that contains it.

A convenient way to produce a subfan is to choose some cones and then
include all of their faces. The intersection condition is inherited
from \(\Sigma\).

\begin{example}
Let \(\Sigma_{\PP^2}\) be the fan of \(\PP^2\), with maximal cones
\[
\sigma_0,\sigma_1,\sigma_2.
\]
The collection consisting of
\[
\sigma_0,\sigma_1
\]
and all of their faces is a subfan. Geometrically, it will correspond
to the union of the two affine charts
\[
V_0\cup V_1
=
\PP^2\setminus\{[0:0:1]\}.
\]
\end{example}

This example previews a general rule:
\[
\text{subfan}
\quad\leadsto\quad
\text{torus-invariant open subvariety}.
\]
We will prove it after constructing \(X_\Sigma\).

\begin{pauseandthink}
Starting with the fan of \(\PP^2\), decide whether each of the
following collections is a subfan:
\begin{enumerate}[label=(\alph*),wide=0pt]
\item one maximal cone and all of its faces;
\item the three rays and the zero cone;
\item the three maximal cones and the zero cone, but no rays;
\item two maximal cones together with only their common ray and the
zero cone.
\end{enumerate}
\end{pauseandthink}

\subsection{Complete fans}

\begin{definition}
A fan \(\Sigma\) in \(N_{\RR}\) is \emph{complete} if
\[
|\Sigma|=N_{\RR}.
\]
\end{definition}

\noindent
\textit{Reference.}
This is Cox--Little--Schenck, Definition~3.1.18(c).

The fans of \(\PP^1\) and \(\PP^2\) are complete: their cones cover
the entire one-parameter subgroup space.

By contrast, the affine fan
\[
\Sigma_\sigma
=
\{\text{faces of }\sigma\}
\]
has support \(\sigma\), so in positive dimension it is not complete.

Completeness has a direct interpretation through one-parameter
subgroups. If \(\Sigma\) is complete, every vector
\[
u\in N
\]
lies in some cone \(\sigma\in\Sigma\). Therefore
\[
\oneps{u}
\]
has a limit in the affine chart \(U_\sigma\). No one-parameter
direction can escape all of the charts.

\begin{prediction}
For \(\PP^1\) and \(\PP^2\), completeness of the fan accompanies
completeness of the variety. Formulate the general relationship you
expect between
\[
|\Sigma|=N_{\RR}
\]
and the geometric property that \(X_\Sigma\) has no points “at
infinity” still missing.

Why might the existence of limits for every one-parameter subgroup be
relevant?
\end{prediction}

\begin{threedictionaries}
The basic fan terminology anticipates global geometry.

\medskip

\begin{center}
\small
\begin{tabularx}{\linewidth}{@{}
  >{\raggedright\arraybackslash}X
  >{\raggedright\arraybackslash}X
  >{\raggedright\arraybackslash}X@{}}
\toprule
\textbf{Algebra}
&
\textbf{Geometry}
&
\textbf{Combinatorics}
\\
\midrule
semigroup algebras for maximal cones
&
an affine open cover
&
\(\Sigma_{\max}\)
\\[1mm]
localizations of those algebras
&
chart intersections
&
common faces
\\[1mm]
restriction to selected charts
&
a torus-invariant open subvariety
&
a subfan
\\[1mm]
every one-parameter subgroup acquires a limit somewhere
&
no direction at infinity is missing
&
\(|\Sigma|=N_{\RR}\)
\\
\bottomrule
\end{tabularx}
\end{center}
\end{threedictionaries}

\begin{guidingquestion}
Given a fan \(\Sigma\), how do the affine varieties \(U_\sigma\) glue
to form a normal separated toric variety \(X_\Sigma\)?
\end{guidingquestion}


\section{Gluing affine toric varieties}
\label{sec:gluing-affine-toric-varieties}
\index{gluing!affine toric varieties}
\index{separation lemma}
\index{toric variety!from a fan}

Let \(\Sigma\) be a fan in \(N_{\RR}\). Every cone
\[
\sigma\in\Sigma
\]
gives a normal affine toric variety
\[
U_\sigma
=
\Spec\Bbbk[S_\sigma],
\qquad
S_\sigma=\sigma^\vee\cap M.
\]
We now glue these affine pieces into one global variety.

The construction is canonical: the fan has already specified every
overlap and every transition map.

\subsection{The canonical overlap of two charts}

Take
\[
\sigma_1,\sigma_2\in\Sigma,
\qquad
\tau=\sigma_1\cap\sigma_2.
\]
By the definition of a fan, \(\tau\) is a face of both cones. Hence
\[
U_{\sigma_1}\supseteq U_\tau\subseteq U_{\sigma_2}
\]
are affine open embeddings.

More explicitly, one can choose
\[
m\in
\sigma_1^\vee\cap(-\sigma_2)^\vee\cap M
\]
such that
\[
\tau
=
\sigma_1\cap m^\perp
=
\sigma_2\cap m^\perp.
\]
Then
\[
U_\tau
=
(U_{\sigma_1})_{\character{m}}
=
(U_{\sigma_2})_{\character{-m}}.
\]

Both copies of \(U_\tau\) have the same coordinate ring
\[
\Bbbk[S_\tau]
\subseteq
\Bbbk[M].
\]
The transition map
\[
g_{\sigma_2,\sigma_1}:
(U_{\sigma_1})_{\character{m}}
\longrightarrow
(U_{\sigma_2})_{\character{-m}}
\]
is therefore induced by the identity map on
\(\Bbbk[S_\tau]\). In particular, it restricts to the identity on the
common dense torus
\[
T=\Spec\Bbbk[M].
\]

\begin{center}
\begin{tikzcd}[column sep=huge,row sep=large]
&
U_\tau
\arrow[dl,hook']
\arrow[dr,hook]
&
\\
U_{\sigma_1}
\arrow[rr,dashed,"\text{glue along }U_\tau"']
&&
U_{\sigma_2}.
\end{tikzcd}
\end{center}

No extra coordinate change has to be invented. It is already encoded
by viewing the same characters in two different affine coordinate
rings.

\subsection{The separation lemma}

The following equality is stronger than merely saying that
\(U_\tau\) is open in both charts.

\begin{proposition}[Separation lemma]
\label{prop:separation-lemma-fans}
Let
\[
\sigma_1,\sigma_2\in\Sigma,
\qquad
\tau=\sigma_1\cap\sigma_2.
\]
Then
\[
S_\tau
=
S_{\sigma_1}+S_{\sigma_2},
\]
where the right-hand side is the Minkowski sum of semigroups.
\end{proposition}

\noindent
\textit{Reference.}
This is Cox--Little--Schenck, Proposition~3.1.3.

\begin{proof}
Since
\[
\sigma_i^\vee\subseteq\tau^\vee,
\]
we immediately have
\[
S_{\sigma_1}+S_{\sigma_2}
\subseteq
S_\tau.
\]

For the reverse inclusion, choose
\[
m\in
\sigma_1^\vee\cap(-\sigma_2)^\vee\cap M
\]
cutting out the common face \(\tau\). Let \(p\in S_\tau\). Because
\(\tau\preceq\sigma_1\), the face-localization equality gives
\[
S_\tau
=
S_{\sigma_1}+\ZZ(-m).
\]
Thus
\[
p=q+\ell(-m)
\]
for some \(q\in S_{\sigma_1}\) and some \(\ell\in\ZZ\).

If \(\ell\geq0\), then
\[
-m\in\sigma_2^\vee\cap M=S_{\sigma_2}
\]
implies
\[
\ell(-m)\in S_{\sigma_2}.
\]
If \(\ell<0\), then
\[
p=q+(-\ell)m\in S_{\sigma_1},
\]
because \(m\in S_{\sigma_1}\). In either case,
\[
p\in S_{\sigma_1}+S_{\sigma_2}.
\]
\end{proof}

The lemma says that every character regular on the overlap can be
written as a product of one character regular on the first chart and
one regular on the second:
\[
\character{p}
=
\character{q}\character{r},
\qquad
q\in S_{\sigma_1},
\quad
r\in S_{\sigma_2}.
\]

This algebraic statement will be the key to separatedness.

\subsection{The cocycle condition is automatic}

Suppose
\[
\sigma_1,\sigma_2,\sigma_3\in\Sigma.
\]
Their triple intersection
\[
\omega
=
\sigma_1\cap\sigma_2\cap\sigma_3
\]
is a common face, and
\[
U_\omega
\]
is an affine open subset of every pairwise overlap.

On \(U_\omega\), all transition maps are induced by the identity on
\[
\Bbbk[S_\omega]\subseteq\Bbbk[M].
\]
Consequently,
\[
g_{\sigma_3,\sigma_2}
\circ
g_{\sigma_2,\sigma_1}
=
g_{\sigma_3,\sigma_1}
\]
wherever all three maps are defined.

Thus the gluing maps satisfy the cocycle condition without an
additional calculation: they are all restrictions of the identity on
the dense torus and its Laurent coordinate ring.

\subsection{Constructing \texorpdfstring{\(X_\Sigma\)}{X Sigma}}

Form the disjoint union
\[
\bigsqcup_{\sigma\in\Sigma}U_\sigma.
\]
Whenever
\[
p\in U_{\sigma_1}
\qquad\text{and}\qquad
q\in U_{\sigma_2}
\]
come from the same point of
\[
U_{\sigma_1\cap\sigma_2},
\]
identify \(p\) and \(q\). The cocycle condition makes this an
equivalence relation.

\begin{definition}
The variety obtained by gluing the affine toric varieties
\[
\{U_\sigma:\sigma\in\Sigma\}
\]
along the canonical common opens
\[
U_{\sigma_1\cap\sigma_2}
\]
is denoted
\[
X_\Sigma.
\]
It is called the \emph{toric variety of the fan \(\Sigma\)}.
\end{definition}

Every \(U_\sigma\) maps isomorphically onto an affine open subset of
\(X_\Sigma\), and
\[
U_{\sigma_1}\cap U_{\sigma_2}
=
U_{\sigma_1\cap\sigma_2}.
\]

Only maximal cones are needed to cover the result:
\[
X_\Sigma
=
\bigcup_{\sigma\in\Sigma_{\max}}U_\sigma.
\]
Indeed, every cone \(\tau\in\Sigma\) is a face of some maximal cone
\(\sigma\), and then
\[
U_\tau\subseteq U_\sigma.
\]

\begin{agcheatsheet}
A variety is \emph{separated} when its diagonal
\[
\Delta_X:X\longrightarrow X\times X,
\qquad
p\longmapsto(p,p),
\]
has closed image.

This is the algebraic-geometric replacement for the Hausdorff
condition. Affine varieties are separated, but an arbitrary gluing of
affine varieties need not be.

To prove that a variety assembled from affine charts is separated, it
is enough to check that the diagonal is closed on every product of two
charts.
\end{agcheatsheet}

\subsection{The toric variety of a fan}

We use the global version of the definition from Chapter~2.

\begin{definition}
A \emph{toric variety} with torus \(T\) is an irreducible variety
\(X\) containing \(T\) as a dense open subset such that the
multiplication action of \(T\) on itself extends to an algebraic
action
\[
T\times X\longrightarrow X.
\]
\end{definition}

\noindent
\textit{Reference.}
This is Cox--Little--Schenck, Definition~3.1.1.

\begin{theorem}[The toric variety of a fan]
\label{thm:toric-variety-of-fan}
Let \(\Sigma\) be a fan in \(N_{\RR}\). Then \(X_\Sigma\) is a normal
separated toric variety with dense torus
\[
T=\Spec\Bbbk[M].
\]
\end{theorem}

\noindent
\textit{Reference.}
This is Cox--Little--Schenck, Theorem~3.1.5. The proof there is written
over \(\CC\); the same algebraic argument works over our base field
\(\Bbbk\).

\begin{proof}[Conceptual proof sketch]
The proof has four parts.

\smallskip

\noindent
\textbf{The torus and its action.}
The zero cone is a face of every cone in \(\Sigma\), so every chart
contains the same open torus
\[
U_{\{0\}}=T.
\]
The gluing maps identify these copies by the identity. The torus
actions on the affine charts agree on every overlap and therefore glue
to an action on \(X_\Sigma\).

\smallskip

\noindent
\textbf{Irreducibility.}
Every chart \(U_\sigma\) is irreducible and contains the same dense
open torus. Their union is therefore irreducible, with \(T\) dense.

\smallskip

\noindent
\textbf{Normality.}
Each \(U_\sigma\) is normal because
\[
S_\sigma=\sigma^\vee\cap M
\]
is saturated. Normality is local, so \(X_\Sigma\) is normal.

\smallskip

\noindent
\textbf{Separatedness.}
For
\[
\tau=\sigma_1\cap\sigma_2,
\]
consider the diagonal on the two charts:
\[
\Delta:
U_\tau
\longrightarrow
U_{\sigma_1}\times U_{\sigma_2}.
\]
On coordinate rings it is induced by
\[
\Delta^*:
\Bbbk[S_{\sigma_1}]
\otimes_\Bbbk
\Bbbk[S_{\sigma_2}]
\longrightarrow
\Bbbk[S_\tau],
\]
\[
\character{m}\otimes\character{n}
\longmapsto
\character{m+n}.
\]
The separation lemma says
\[
S_\tau=S_{\sigma_1}+S_{\sigma_2},
\]
so \(\Delta^*\) is surjective. Hence \(\Delta(U_\tau)\) is closed in
\[
U_{\sigma_1}\times U_{\sigma_2}.
\]
Checking this for every pair of affine charts proves that
\(X_\Sigma\) is separated.
\end{proof}

\begin{warningbox}
The affine pieces alone do not guarantee separatedness.

If two copies of \(\AAA^1\) are glued along
\[
\Gm=\AAA^1\setminus\{0\}
\]
by the identity, the result has two distinct origins that cannot be
separated by disjoint open sets. This is the \emph{affine line with
doubled origin}, and it is not separated.

The fan construction rules out this pathology. Its canonical
overlaps satisfy the separation lemma, which forces the diagonal to
be closed.
\end{warningbox}


\begin{example}[The toric blowup of \texorpdfstring{\(\AAA^2\)}{A2} at the origin]
\label{ex:blowup-a2-origin-fan}
Start with the affine fan of the first quadrant
\[
\Sigma_{\AAA^2}
=
\bigl\{
\{0\},\rho_1,\rho_2,\sigma
\bigr\},
\qquad
\rho_1=\Cone{e_1},\ \rho_2=\Cone{e_2},\ 
\sigma=\Cone{e_1,e_2}.
\]
Insert the new ray
\[
\rho=\Cone{e_1+e_2}.
\]
This subdivides \(\sigma\) into two smooth cones
\[
\sigma_1=\Cone{e_1,e_1+e_2},
\qquad
\sigma_2=\Cone{e_1+e_2,e_2}.
\]

\begin{figure}[ht]
\centering
\begin{tikzpicture}[scale=1.08]
  % left: original fan
  \begin{scope}[xshift=-4.5cm]
    \fill[PastelYellow] (0,0)--(2.5,0)--(0,2.5)--cycle;
    \draw[->,black!50] (-0.35,0)--(3.0,0);
    \draw[->,black!50] (0,-0.35)--(0,3.0);
    \draw[very thick,DictionaryBlue,-{Latex[length=2.3mm]}] (0,0)--(2.7,0);
    \draw[very thick,SurfaceGreen,-{Latex[length=2.3mm]}] (0,0)--(0,2.7);
    \node[DictionaryBlue,below] at (2.55,0) {\(\rho_1\)};
    \node[SurfaceGreen,left] at (0,2.55) {\(\rho_2\)};
    \node at (0.8,0.8) {\(\sigma\)};
    \node[font=\bfseries] at (0.8,3.25) {fan of \(\AAA^2\)};
  \end{scope}

  % arrow
  \draw[very thick,-{Latex[length=2.6mm]}] (-0.75,0.8)--(0.75,0.8);
  \node[align=center] at (0,1.25) {star\\ subdivision};

  % right: blowup fan
  \begin{scope}[xshift=4.5cm]
    \fill[DictionaryBlue!10] (0,0)--(2.5,0)--(1.25,1.25)--cycle;
    \fill[SurfaceGreen!10] (0,0)--(1.25,1.25)--(0,2.5)--cycle;
    \draw[->,black!50] (-0.35,0)--(3.0,0);
    \draw[->,black!50] (0,-0.35)--(0,3.0);
    \draw[very thick,DictionaryBlue,-{Latex[length=2.3mm]}] (0,0)--(2.7,0);
    \draw[very thick,SurfaceGreen,-{Latex[length=2.3mm]}] (0,0)--(0,2.7);
    \draw[very thick,PredictionTeal,-{Latex[length=2.3mm]}] (0,0)--(2.15,2.15);
    \node[DictionaryBlue,below] at (2.55,0) {\(\rho_1\)};
    \node[SurfaceGreen,left] at (0,2.55) {\(\rho_2\)};
    \node[PredictionTeal,above right] at (1.95,1.95) {\(\rho\)};
    \node at (1.55,0.45) {\(\sigma_1\)};
    \node at (0.45,1.55) {\(\sigma_2\)};
    \node[font=\bfseries] at (1.0,3.25)
      {fan of \(\Bl_0(\AAA^2)\)};
  \end{scope}
\end{tikzpicture}
\caption{Subdividing the first quadrant by the ray through
\(e_1+e_2\) produces the fan of the toric blowup of \(\AAA^2\) at the
origin.}
\label{fig:fan-blowup-a2-origin}
\end{figure}

Each maximal cone is smooth, so each affine chart is isomorphic to
\(\AAA^2\). Indeed,
\[
U_{\sigma_1}\simeq\Spec\Bbbk[x,xy^{-1}],
\qquad
U_{\sigma_2}\simeq\Spec\Bbbk[y,x^{-1}y].
\]
The identity map of lattices
\[
N\longrightarrow N
\]
is compatible from the subdivided fan to the single-cone fan of
\(\AAA^2\), so it induces a toric morphism
\[
\pi:X_{\Sigma'}\longrightarrow \AAA^2.
\]

Geometrically, \(\pi\) is the blowup of the origin. The new ray
\(\rho\) contributes a new one-dimensional orbit; its closure is the
exceptional divisor
\[
V(\rho)\simeq\PP^1.
\]
The old rays still correspond to the strict transforms of the two
coordinate axes.
\end{example}

\begin{threedictionaries}
The gluing construction is the global form of the three dictionaries.

\medskip

\begin{center}
\small
\begin{tabularx}{\linewidth}{@{}
  >{\raggedright\arraybackslash}X
  >{\raggedright\arraybackslash}X
  >{\raggedright\arraybackslash}X@{}}
\toprule
\textbf{Algebra}
&
\textbf{Geometry}
&
\textbf{Combinatorics}
\\
\midrule
\(\Bbbk[S_\sigma]\)
&
an affine chart \(U_\sigma\)
&
a cone \(\sigma\)
\\[1mm]
localization to
\(\Bbbk[S_{\sigma_1\cap\sigma_2}]\)
&
the overlap
\(U_{\sigma_1}\cap U_{\sigma_2}\)
&
the common face
\(\sigma_1\cap\sigma_2\)
\\[1mm]
identity maps inside \(\Bbbk[M]\)
&
transition maps satisfying the cocycle condition
&
fan compatibility
\\[1mm]
surjectivity from
\(S_{\sigma_1\cap\sigma_2}
=S_{\sigma_1}+S_{\sigma_2}\)
&
separatedness
&
the separation lemma
\\
\bottomrule
\end{tabularx}
\end{center}
\end{threedictionaries}

\begin{guidingquestion}
Once the affine pieces have been glued, how are the torus orbits of
\(X_\Sigma\) indexed by the cones of the fan, and how do face
relations control orbit closures?
\end{guidingquestion}


\section{The orbit--cone correspondence}
\label{sec:orbit-cone-correspondence}
\index{orbit--cone correspondence}
\index{torus orbit}
\index{orbit closure}
\index{star of a cone}

A fan is more than an atlas of affine charts. Its individual cones
also index the torus orbits of the glued variety.

The starting point is the distinguished limit attached to a cone.

\subsection{Distinguished points}

Let
\[
\sigma\in\Sigma,
\qquad
S_\sigma=\sigma^\vee\cap M.
\]
Define a point
\[
\gamma_\sigma\in U_\sigma
\]
by the semigroup homomorphism
\[
\gamma_\sigma(m)
=
\begin{cases}
1,&m\in\sigma^\perp\cap M,\\
0,&m\notin\sigma^\perp.
\end{cases}
\]

This is the \emph{distinguished point} corresponding to \(\sigma\).

If
\[
u\in\relint(\sigma)\cap N,
\]
then the one-parameter subgroup criterion from Chapter~2 gives
\[
\lim_{t\to0}\oneps{u}(t)=\gamma_\sigma.
\]
Thus every lattice direction in the relative interior of \(\sigma\)
selects the same distinguished point.

\noindent
\textit{Reference.}
The distinguished point and its interpretation by one-parameter
limits are Cox--Little--Schenck, Proposition~3.2.2 and the discussion
preceding it.

Now let the torus move this point.

\begin{definition}
For \(\sigma\in\Sigma\), define
\[
O(\sigma)
=
T\cdot\gamma_\sigma
\subseteq X_\Sigma.
\]
\end{definition}

\subsection{What kind of torus is \texorpdfstring{\(O(\sigma)\)}{O(sigma)}?}

Let
\[
N_\sigma
=
N\cap\Span_{\RR}(\sigma)
\]
and set
\[
N(\sigma)=N/N_\sigma.
\]
The lattice dual to \(N(\sigma)\) is
\[
\sigma^\perp\cap M.
\]

\begin{proposition}[The orbit attached to a cone]
\label{prop:orbit-attached-to-cone}
There are natural isomorphisms
\[
O(\sigma)
\simeq
\Hom_{\ZZ}
\bigl(\sigma^\perp\cap M,\Bbbk^\times\bigr)
\simeq
T_{N(\sigma)}.
\]
Consequently,
\[
\dim O(\sigma)
=
\rank N-\dim\sigma.
\]
\end{proposition}

\noindent
\textit{Reference.}
This is Cox--Little--Schenck, Lemmas~3.2.4 and~3.2.5.

\begin{proof}[Conceptual proof sketch]
A point of \(U_\sigma\) is a semigroup homomorphism
\[
\gamma:S_\sigma\longrightarrow\Bbbk.
\]
The orbit of \(\gamma_\sigma\) consists exactly of those \(\gamma\)
satisfying
\[
\gamma(m)\neq0
\quad\Longleftrightarrow\quad
m\in\sigma^\perp\cap M.
\]
Restricting such a \(\gamma\) to
\(\sigma^\perp\cap M\) gives a group homomorphism to
\(\Bbbk^\times\), and every such group homomorphism extends by zero
outside \(\sigma^\perp\).

The perfect pairing
\[
(\sigma^\perp\cap M)\times N(\sigma)\longrightarrow\ZZ
\]
identifies this group of homomorphisms with the torus \(T_{N(\sigma)}\).
Its dimension is
\[
\rank N(\sigma)
=
\rank N-\rank N_\sigma
=
\rank N-\dim\sigma.
\]
\end{proof}

Two extreme cases are already familiar:
\[
O(\{0\})=T,
\]
while a full-dimensional cone gives a zero-dimensional orbit, hence a
torus-fixed point.

\begin{pauseandthink}
Let \(\Sigma\) be a fan in a rank-\(n\) lattice.

What orbit dimensions correspond to:
\[
\text{the zero cone},\qquad
\text{a ray},\qquad
\text{a maximal \(n\)-dimensional cone}?
\]
What geometric strata should the rays represent?
\end{pauseandthink}

\subsection{The theorem}

\begin{theorem}[Orbit--cone correspondence]
\label{thm:orbit-cone-correspondence}
Let \(X_\Sigma\) be the toric variety of a fan
\(\Sigma\subseteq N_{\RR}\). Then:

\begin{enumerate}[label=(\alph*),wide=0pt]
\item The assignment
\[
\sigma\longmapsto O(\sigma)
\]
is a bijection
\[
\{\text{cones of }\Sigma\}
\longleftrightarrow
\{\text{\(T\)-orbits in }X_\Sigma\}.
\]

\item For every cone,
\[
\dim O(\sigma)
=
\rank N-\dim\sigma.
\]

\item Each affine chart decomposes as
\[
U_\sigma
=
\bigsqcup_{\tau\preceq\sigma}O(\tau).
\]

\item Face relations reverse under orbit closure:
\[
\tau\preceq\sigma
\quad\Longleftrightarrow\quad
O(\sigma)\subseteq\overline{O(\tau)}.
\]
In particular,
\[
\overline{O(\tau)}
=
\bigsqcup_{\tau\preceq\sigma}O(\sigma).
\]
\end{enumerate}
\end{theorem}

\noindent
\textit{Reference.}
This is Cox--Little--Schenck, Theorem~3.2.6.

\begin{proof}[Conceptual proof sketch]
The previous proposition constructs one orbit \(O(\sigma)\) for every
cone.

Conversely, let \(O\) be any orbit. Among the cones whose affine
charts meet \(O\), choose a minimal one, say \(\sigma\). Represent a
point of \(O\cap U_\sigma\) by a semigroup homomorphism
\[
\gamma:S_\sigma\to\Bbbk.
\]
The set of exponents on which \(\gamma\) is nonzero is the lattice
part of a face of \(\sigma^\vee\). Duality identifies that face with
a face \(\tau\preceq\sigma\), and then the point lies in \(U_\tau\).
Minimality forces
\[
\tau=\sigma,
\]
so the orbit is \(O(\sigma)\). This proves the bijection.

The dimension formula follows from
\[
O(\sigma)\simeq T_{N(\sigma)}.
\]
The description of \(U_\sigma\) follows because its orbits are exactly
those indexed by faces of \(\sigma\).

Finally, if \(\tau\preceq\sigma\), move the distinguished point
\(\gamma_\tau\) by a one-parameter subgroup whose lattice vector lies
in \(\relint(\sigma)\). Its limit lies in
\(O(\sigma)\), showing
\[
O(\sigma)\subseteq\overline{O(\tau)}.
\]
The converse follows from the affine-chart decomposition. The
technical part of the complete proof verifies that the displayed
union is Zariski closed.
\end{proof}

The theorem reverses both dimension and incidence:
\[
\begin{array}{c@{\qquad\longleftrightarrow\qquad}c}
\text{larger cone}
&
\text{smaller orbit}
\\[1mm]
\tau\preceq\sigma
&
O(\sigma)\subseteq\overline{O(\tau)}.
\end{array}
\]

\begin{corollary}
\label{cor:closed-orbit-affine-chart}
For every cone \(\sigma\in\Sigma\), the affine chart \(U_\sigma\)
contains a unique closed torus orbit, namely \(O(\sigma)\).
\end{corollary}

\begin{proof}
The chart contains precisely the orbits \(O(\tau)\) with
\(\tau\preceq\sigma\). The closure of every such orbit contains
\(O(\sigma)\), while no strictly larger cone belongs to the affine fan
of \(\sigma\).
\end{proof}


\subsection{Reading the blowup fan}
\label{subsec:reading-blowup-fan}

Return to the blowup fan from \cref{ex:blowup-a2-origin-fan}. It has
three rays
\[
\rho_1=\Cone{e_1},
\qquad
\rho=\Cone{e_1+e_2},
\qquad
\rho_2=\Cone{e_2},
\]
and two maximal cones
\[
\sigma_1=\Cone{e_1,e_1+e_2},
\qquad
\sigma_2=\Cone{e_1+e_2,e_2}.
\]

The orbit--cone correspondence can be read directly from the picture.

\begin{figure}[ht]
\centering
\begin{tikzpicture}[scale=1.0,
  orbitbox/.style={rounded corners, draw=black!50, inner sep=4pt, align=center},
  hasse/.style={circle, draw=black!50, inner sep=1.5pt, minimum size=7mm}
]
  % Fan picture
  \begin{scope}[xshift=-4.8cm]
    \fill[DictionaryBlue!12] (0,0)--(2.6,0)--(1.3,1.3)--cycle;
    \fill[SurfaceGreen!12] (0,0)--(1.3,1.3)--(0,2.6)--cycle;
    \draw[->,black!50] (-0.35,0)--(3.0,0);
    \draw[->,black!50] (0,-0.35)--(0,3.0);
    \draw[very thick,DictionaryBlue,-{Latex[length=2.3mm]}] (0,0)--(2.8,0);
    \draw[very thick,PredictionTeal,-{Latex[length=2.3mm]}] (0,0)--(2.2,2.2);
    \draw[very thick,SurfaceGreen,-{Latex[length=2.3mm]}] (0,0)--(0,2.8);

    \node[DictionaryBlue,below] at (2.62,0) {\(\rho_1\)};
    \node[PredictionTeal,above right] at (1.98,1.98) {\(\rho\)};
    \node[SurfaceGreen,left] at (0,2.62) {\(\rho_2\)};

    \node[fill=DictionaryBlue!18,rounded corners,inner sep=2pt] at (1.7,0.45) {\(\sigma_1\)};
    \node[fill=SurfaceGreen!18,rounded corners,inner sep=2pt] at (0.45,1.7) {\(\sigma_2\)};
    \node[fill=white,rounded corners,inner sep=2pt] at (-0.38,-0.22) {\(\{0\}\)};
    \node[font=\bfseries] at (0.95,3.3) {cones in the fan};
  \end{scope}

  % Bijection symbol
  \node[font=\Large] at (0,0.55) {\(\longleftrightarrow\)};

  % Hasse diagram for orbit closures
  \begin{scope}[xshift=4.9cm]
    \node[font=\bfseries] at (0,3.3) {orbit-closure poset};

    % nodes
    \node[hasse, fill=white] (T) at (0,2.5) {\(\Gm^2\)};
    \node[hasse, fill=DictionaryBlue!12] (R1) at (-2.0,1.1) {\(O(\rho_1)\)};
    \node[hasse, fill=PredictionTeal!12] (R) at (0,1.1) {\(O(\rho)\)};
    \node[hasse, fill=SurfaceGreen!12] (R2) at (2.0,1.1) {\(O(\rho_2)\)};
    \node[hasse, fill=DictionaryBlue!18] (S1) at (-1.0,-0.3) {\(O(\sigma_1)\)};
    \node[hasse, fill=SurfaceGreen!18] (S2) at (1.0,-0.3) {\(O(\sigma_2)\)};

    % edges showing covering relations in closure order
    \draw[black!65] (T)--(R1);
    \draw[black!65] (T)--(R);
    \draw[black!65] (T)--(R2);
    \draw[black!65] (R1)--(S1);
    \draw[black!65] (R)--(S1);
    \draw[black!65] (R)--(S2);
    \draw[black!65] (R2)--(S2);

    % side labels
    
  \end{scope}
\end{tikzpicture}
\caption{Color-coded orbit--cone correspondence for the fan of
\(\Bl_0(\AAA^2)\). The colors and labels establish the bijection
between cones and orbits, while the Hasse diagram records the
orbit-closure order. Ordering cones by inclusion and orbit closures by
reverse inclusion, the correspondence is an isomorphism of posets.}
\label{fig:blowup-orbit-cone-bijection}
\end{figure}

The figure makes three features visible at once.

\begin{itemize}
\item The zero cone corresponds to the dense torus.
\item The three rays correspond to the three boundary curves: the two
strict transforms of the coordinate axes and the exceptional divisor.
\item The two maximal cones correspond to the two torus-fixed points,
which lie on the exceptional divisor.
\end{itemize}

This is a good first reminder that the orbit decomposition is not only
a list of strata: it also remembers how those strata fit together.

\subsection{Reading the fan of \texorpdfstring{\(\PP^2\)}{P2}}

Let
\[
\rho_i=\Cone{e_i},
\qquad
i=0,1,2,
\]
and let \(\sigma_0,\sigma_1,\sigma_2\) be the three maximal cones from
Section~3.1. The complete orbit decomposition is:

\begin{center}
\small
\renewcommand{\arraystretch}{1.25}
\begin{tabularx}{\linewidth}{@{}
  >{\raggedright\arraybackslash}p{0.19\linewidth}
  >{\raggedright\arraybackslash}p{0.13\linewidth}
  >{\raggedright\arraybackslash}X
  >{\raggedright\arraybackslash}p{0.23\linewidth}@{}}
\toprule
\textbf{Cones}
&
\textbf{Orbit dimension}
&
\textbf{Orbit in \(\PP^2\)}
&
\textbf{Orbit closure}
\\
\midrule
\(\{0\}\)
&
\(2\)
&
\(\{[z_0:z_1:z_2]:z_0z_1z_2\neq0\}\simeq\Gm^2\)
&
\(\PP^2\)
\\
\(\rho_0,\rho_1,\rho_2\)
&
\(1\)
&
a coordinate line with its two coordinate points removed
&
a coordinate line \(\simeq\PP^1\)
\\
\(\sigma_0,\sigma_1,\sigma_2\)
&
\(0\)
&
the three coordinate points
&
the point itself
\\
\bottomrule
\end{tabularx}
\end{center}

\begin{surfacelaboratory}
Use the fan and the table to answer the following.

\begin{enumerate}[label=(\alph*),wide=0pt]
\item Why does each ray belong to exactly two maximal cones, and how
does this predict that the closure of the corresponding
one-dimensional orbit contains exactly two fixed points?

\item Choose one maximal cone \(\sigma_i\). List its faces and recover
the three orbit types appearing in its affine-plane chart.

\item Starting from a coordinate line in \(\PP^2\), identify its open
torus orbit and the two lower-dimensional orbits in its boundary.

\item Count the orbits of each dimension directly from
\[
\Sigma(0),\qquad\Sigma(1),\qquad\Sigma(2).
\]
\end{enumerate}
\end{surfacelaboratory}

\subsection{Orbit closures are toric varieties}

For \(\tau\in\Sigma\), write
\[
V(\tau)=\overline{O(\tau)}.
\]
The theorem gives
\[
V(\tau)
=
\bigsqcup_{\tau\preceq\sigma}O(\sigma).
\]

To describe \(V(\tau)\) by a fan, project every cone containing
\(\tau\) to
\[
N(\tau)_{\RR}
=
N_{\RR}/\Span_{\RR}(\tau).
\]
If \(\overline{\sigma}\) denotes the image of \(\sigma\), define
\[
\Star(\tau)
=
\left\{
\overline{\sigma}\subseteq N(\tau)_{\RR}:
\tau\preceq\sigma\in\Sigma
\right\}.
\]

\begin{proposition}[The fan of an orbit closure]
\label{prop:fan-of-orbit-closure}
The collection \(\Star(\tau)\) is a fan, and
\[
V(\tau)
\simeq
X_{\Star(\tau)}.
\]
\end{proposition}

\noindent
\textit{Reference.}
This is Cox--Little--Schenck, Proposition~3.2.7.

\begin{proof}[Conceptual proof sketch]
The orbit \(O(\tau)\) is the torus with one-parameter subgroup lattice
\[
N(\tau)=N/N_\tau.
\]
The orbits contained in its closure are indexed exactly by the cones
\(\sigma\) containing \(\tau\). After quotienting by the directions of
\(\tau\), these cones form the fan \(\Star(\tau)\).
The orbit--cone correspondence for that quotient fan reproduces the
same orbits and the same incidence relations as \(V(\tau)\).
\end{proof}

For a ray \(\rho\) in a surface fan,
\[
N(\rho)
\]
has rank one. Hence \(V(\rho)\) is a toric curve. In the fan of
\(\PP^2\), the two maximal cones containing \(\rho\) project to the two
opposite rays in \(N(\rho)_{\RR}\), so
\[
V(\rho)\simeq\PP^1.
\]

\begin{threedictionaries}
The orbit--cone correspondence translates the face structure of the
fan into the stratification of the variety.

\medskip

\begin{center}
\small
\begin{tabularx}{\linewidth}{@{}
  >{\raggedright\arraybackslash}X
  >{\raggedright\arraybackslash}X
  >{\raggedright\arraybackslash}X@{}}
\toprule
\textbf{Algebra}
&
\textbf{Geometry}
&
\textbf{Combinatorics}
\\
\midrule
characters in \(\sigma^\perp\cap M\)
&
the torus \(O(\sigma)\)
&
a cone \(\sigma\)
\\[1mm]
rank of \(\sigma^\perp\cap M\)
&
\(\dim O(\sigma)\)
&
\(\rank N-\dim\sigma\)
\\[1mm]
vanishing of additional characters
&
movement toward the boundary
&
passing to a larger cone
\\[1mm]
quotient character and one-parameter lattices
&
the orbit closure \(V(\tau)\)
&
the fan \(\Star(\tau)\)
\\
\bottomrule
\end{tabularx}
\end{center}
\end{threedictionaries}

\begin{guidingquestion}
Which global geometric properties of \(X_\Sigma\)—smoothness,
completeness, products, and morphisms—can be recognized directly from
the fan?
\end{guidingquestion}


\section{Properties read from the fan}
\label{sec:properties-read-from-fan}
\index{smoothness criterion!fan}
\index{completeness!fan criterion}
\index{product fan}
\index{toric morphism!fan criterion}

The construction
\[
\Sigma\longmapsto X_\Sigma
\]
turns combinatorial conditions into geometric properties. The most
basic entries are:

\begin{center}
\small
\renewcommand{\arraystretch}{1.25}
\begin{tabularx}{0.92\linewidth}{@{}
  >{\raggedright\arraybackslash}X
  >{\raggedright\arraybackslash}X@{}}
\toprule
\textbf{Geometry of \(X_\Sigma\)}
&
\textbf{Combinatorics of \(\Sigma\)}
\\
\midrule
smooth
&
every cone is smooth
\\
finite quotient singularities
&
every cone is simplicial
\\
complete
&
\(|\Sigma|=N_{\RR}\)
\\
product
&
product fan
\\
toric morphism
&
a compatible lattice map
\\
\bottomrule
\end{tabularx}
\end{center}

We now explain these correspondences, concentrating on the criterion
and the idea behind it.

\subsection{Smoothness is checked cone by cone}

\begin{definition}
A fan \(\Sigma\) is \emph{smooth} if every cone
\(\sigma\in\Sigma\) is smooth: the primitive generators of its rays
form part of a basis of \(N\).

A fan is \emph{simplicial} if every cone is generated by linearly
independent vectors.
\end{definition}

In dimension two, every strongly convex cone is simplicial, but it
need not be smooth. If
\[
\sigma=\Cone{u_1,u_2}
\]
has primitive ray generators, then
\[
\sigma\text{ is smooth}
\quad\Longleftrightarrow\quad
|\det(u_1,u_2)|=1.
\]

\begin{theorem}[Global smoothness criterion]
\label{thm:global-smoothness-criterion}
The toric variety \(X_\Sigma\) is smooth if and only if the fan
\(\Sigma\) is smooth.
\end{theorem}

\noindent
\textit{Reference.}
This is Cox--Little--Schenck, Theorem~3.1.19(a).

\begin{proof}
The affine charts \(U_\sigma\), for
\(\sigma\in\Sigma_{\max}\), cover \(X_\Sigma\). Smoothness is local,
and the affine smoothness criterion gives
\[
U_\sigma\text{ smooth}
\quad\Longleftrightarrow\quad
\sigma\text{ smooth}.
\]
Thus every point of \(X_\Sigma\) is smooth exactly when every maximal
cone is smooth. Since every cone is a face of a maximal cone, this is
equivalent to smoothness of the entire fan.
\end{proof}

\begin{observation}
Over a field of characteristic zero, one also has
\[
X_\Sigma\text{ has only finite quotient singularities}
\quad\Longleftrightarrow\quad
\Sigma\text{ is simplicial}.
\]
This is Cox--Little--Schenck, Theorem~3.1.19(b). Its proof requires
more singularity theory than we need here and will be revisited later.

The cyclic quotient surface from the challenge exercise of Chapter~2
is the local model: the cone is simplicial, but its primitive
generators need not form a basis of the chosen lattice.
\end{observation}

\begin{pauseandthink}
Let \(\Sigma\) be a two-dimensional fan. Explain why its singular
points can occur only at the zero-dimensional orbits associated to
nonsmooth maximal cones.

Why can the one-dimensional orbits not be singular?
\end{pauseandthink}

\subsection{Completeness and one-parameter limits}

Recall that a fan is complete when
\[
|\Sigma|=N_{\RR}.
\]
This condition has an exact geometric meaning.

\begin{agcheatsheet}
A variety \(X\) is \emph{complete} if its structure morphism
\[
X\longrightarrow\Spec\Bbbk
\]
is proper.

Over \(\CC\), complete varieties are precisely those whose complex
points are compact in the classical topology. Every projective variety
is complete, although a complete variety need not be projective.
\end{agcheatsheet}

\begin{theorem}[Completeness criterion]
\label{thm:completeness-criterion-fan}
For a toric variety \(X_\Sigma\), the following are equivalent:
\begin{enumerate}[label=(\alph*),wide=0pt]
\item \(X_\Sigma\) is complete;
\item every one-parameter subgroup
\[
\oneps{u}:\Gm\longrightarrow T,
\qquad
u\in N,
\]
has a limit in \(X_\Sigma\) as \(t\to0\);
\item the fan is complete:
\[
|\Sigma|=N_{\RR}.
\]
\end{enumerate}
\end{theorem}

\noindent
\textit{Reference.}
This is Cox--Little--Schenck, Theorem~3.4.6.

\begin{proof}[Conceptual proof sketch]
The equivalence of (b) and (c) is the one-parameter subgroup
criterion applied chart by chart:
\[
\lim_{t\to0}\oneps{u}(t)\text{ exists in }U_\sigma
\quad\Longleftrightarrow\quad
u\in\sigma.
\]
Therefore the limit exists somewhere in \(X_\Sigma\) exactly when
\[
u\in|\Sigma|.
\]

If the support is not all of \(N_{\RR}\), rationality and finiteness
of the fan allow us to choose a lattice vector
\[
u\notin|\Sigma|.
\]
The corresponding map \(\oneps{u}\) has no extension across \(t=0\),
so the valuative criterion rules out completeness.

Conversely, if the support is complete, every one-parameter subgroup
has a limit in some affine chart. The substantive part of the proof
uses the valuative criterion to show that, for toric varieties, these
torus limits suffice to establish properness.
\end{proof}

This criterion makes the pictures of \(\AAA^2\) and \(\PP^2\)
immediately different. The affine-plane fan occupies only one
quadrant, while the three maximal cones of the projective-plane fan
cover all of \(N_{\RR}\).

\begin{prediction}
A complete fan in \(\RR^2\) must surround the origin. What is the
smallest possible number of rays in a complete two-dimensional fan?
What toric surface should the smallest smooth example produce?
\end{prediction}

\subsection{Product fans}

Let \(\Sigma_1\) and \(\Sigma_2\) be fans in
\((N_1)_{\RR}\) and \((N_2)_{\RR}\).

\begin{definition}
Their \emph{product fan} is
\[
\Sigma_1\times\Sigma_2
=
\left\{
\sigma_1\times\sigma_2:
\sigma_i\in\Sigma_i
\right\}
\]
in
\[
(N_1\oplus N_2)_{\RR}.
\]
\end{definition}

Faces and intersections are computed factor by factor, so this is
indeed a fan.

\begin{proposition}[Products]
\label{prop:product-fans}
There is a canonical isomorphism
\[
X_{\Sigma_1\times\Sigma_2}
\simeq
X_{\Sigma_1}\times X_{\Sigma_2}.
\]
\end{proposition}

\noindent
\textit{Reference.}
This is Cox--Little--Schenck, Proposition~3.1.14.

\begin{proof}
For a product cone,
\[
(\sigma_1\times\sigma_2)^\vee
=
\sigma_1^\vee\times\sigma_2^\vee,
\]
and hence
\[
S_{\sigma_1\times\sigma_2}
=
S_{\sigma_1}\oplus S_{\sigma_2}.
\]
Therefore
\[
\Bbbk[S_{\sigma_1\times\sigma_2}]
\simeq
\Bbbk[S_{\sigma_1}]
\otimes_\Bbbk
\Bbbk[S_{\sigma_2}],
\]
which gives
\[
U_{\sigma_1\times\sigma_2}
\simeq
U_{\sigma_1}\times U_{\sigma_2}.
\]
These affine isomorphisms agree on common face charts and therefore
glue.
\end{proof}

Smoothness and completeness behave as expected:
\[
\Sigma_1\times\Sigma_2\text{ is smooth}
\quad\Longleftrightarrow\quad
\Sigma_1,\Sigma_2\text{ are smooth},
\]
and
\[
|\Sigma_1\times\Sigma_2|
=
|\Sigma_1|\times|\Sigma_2|.
\]

\subsection{Toric morphisms from lattice maps}

Let \(\Sigma_i\) be a fan in \((N_i)_{\RR}\), and let
\[
F:N_1\longrightarrow N_2
\]
be a lattice homomorphism.

\begin{definition}
The map \(F\) is \emph{compatible with the fans}
\(\Sigma_1,\Sigma_2\) if, for every cone
\(\sigma\in\Sigma_1\), there is a cone
\(\sigma'\in\Sigma_2\) such that
\[
F_{\RR}(\sigma)\subseteq\sigma'.
\]
\end{definition}

This is the global version of the affine cone criterion.

\begin{theorem}[Toric morphisms and maps of fans]
\label{thm:toric-morphisms-maps-of-fans}
Compatible lattice maps
\[
F:(N_1,\Sigma_1)\longrightarrow(N_2,\Sigma_2)
\]
are in natural correspondence with toric morphisms
\[
\phi_F:X_{\Sigma_1}\longrightarrow X_{\Sigma_2}.
\]
\end{theorem}

\noindent
\textit{Reference.}
The definition of toric morphism is Cox--Little--Schenck,
Definition~3.3.3, and the correspondence is Theorem~3.3.4.

\begin{proof}[Conceptual proof sketch]
Suppose \(F\) is compatible. For every
\(\sigma\in\Sigma_1\), choose
\(\sigma'\in\Sigma_2\) with
\[
F_{\RR}(\sigma)\subseteq\sigma'.
\]
The affine cone criterion produces a toric morphism
\[
U_\sigma\longrightarrow U_{\sigma'}.
\]
All of these maps restrict to the same torus homomorphism, so they
agree on common face charts and glue to a global morphism.

Conversely, a toric morphism restricts on the dense tori to a group
homomorphism, hence induces a lattice map
\[
F:N_1\to N_2.
\]
The image of each affine chart \(U_\sigma\) meets a suitable affine
chart of the target. Applying the affine cone criterion shows that
\(F_{\RR}(\sigma)\) lies in a cone of \(\Sigma_2\).

The bookkeeping needed to choose target charts and verify gluing is
the only global step.
\end{proof}

The product projections are the basic examples. The lattice map
\[
\operatorname{pr}_1:
N_1\oplus N_2\longrightarrow N_1
\]
sends
\[
\sigma_1\times\sigma_2
\longmapsto
\sigma_1,
\]
so it induces
\[
X_{\Sigma_1}\times X_{\Sigma_2}
\longrightarrow
X_{\Sigma_1}.
\]

\begin{observation}
Finite-index changes of lattice fit the same theorem. If
\[
N'\subseteq N
\]
has finite index and the same real fan \(\Sigma\) is viewed in both
lattices, then the inclusion is compatible and induces
\[
X_{\Sigma,N'}
\longrightarrow
X_{\Sigma,N}.
\]
This is the global version of the lattice change behind the cyclic
quotient orbifolds in the Chapter~2 challenge.
\end{observation}

\begin{surfacelaboratory}
Let \(\Sigma_{\square}\) be the fan in \(\RR^2\) whose rays are
\[
\RR_{\geq0}e_1,\quad
\RR_{\geq0}e_2,\quad
\RR_{\geq0}(-e_1),\quad
\RR_{\geq0}(-e_2),
\]
and whose maximal cones are the four quadrants.

\begin{enumerate}[label=(\alph*),wide=0pt]
\item Verify directly that \(\Sigma_{\square}\) is smooth and
complete.

\item Show that it is the product of two copies of the fan of
\(\PP^1\).

\item Deduce that
\[
X_{\Sigma_{\square}}
\simeq
\PP^1\times\PP^1.
\]

\item Use the coordinate projections
\[
\ZZ^2\longrightarrow\ZZ
\]
to recover the two morphisms
\[
\PP^1\times\PP^1\longrightarrow\PP^1.
\]

\item Count the two-dimensional, one-dimensional, and zero-dimensional
torus orbits directly from the fan.
\end{enumerate}
\end{surfacelaboratory}

\begin{guidingquestion}
How do the standard fans of projective spaces, products, weighted
projective surfaces, and blowups display these properties in concrete
pictures?
\end{guidingquestion}


\section{A catalogue of fans and toric surfaces}
\label{sec:catalogue-fans-toric-surfaces}
\index{projective plane!fan}
\index{weighted projective plane}
\index{Hirzebruch surface}
\index{blowup!star subdivision}

We finish the chapter by collecting a small number of surface fans
that will reappear throughout the course. The goal is not to memorize
the pictures, but to learn which features should be read immediately:

\[
\text{support},\qquad
\text{primitive ray generators},\qquad
\text{determinants of adjacent rays},\qquad
\text{subdivisions}.
\]

\subsection{Four recurring complete fans}

The following four fans form a useful visual vocabulary.

\begin{figure}[ht]
\centering
\begin{tikzpicture}[scale=0.78]
  % P2
  \begin{scope}[xshift=-5.4cm,yshift=3.7cm]
    \fill[DictionaryBlue!9] (0,0)--(2.1,0)--(0,2.1)--cycle;
    \fill[SurfaceGreen!9] (0,0)--(0,2.1)--(-1.55,-1.55)--cycle;
    \fill[ExampleGold!12] (0,0)--(-1.55,-1.55)--(2.1,0)--cycle;
    \draw[very thick,DictionaryBlue,-{Latex[length=2.2mm]}] (0,0)--(2.4,0);
    \draw[very thick,SurfaceGreen,-{Latex[length=2.2mm]}] (0,0)--(0,2.4);
    \draw[very thick,ExampleGold,-{Latex[length=2.2mm]}] (0,0)--(-1.8,-1.8);
    \node[font=\bfseries] at (0.35,2.95) {\(\PP^2\)};
    \node[below] at (2.2,0) {\((1,0)\)};
    \node[left] at (0,2.2) {\((0,1)\)};
    \node[above left] at (-1.65,-1.65) {\((-1,-1)\)};
  \end{scope}

  % P1xP1
  \begin{scope}[xshift=4.8cm,yshift=3.7cm]
    \fill[DictionaryBlue!8] (0,0)--(2.1,0)--(0,2.1)--cycle;
    \fill[SurfaceGreen!8] (0,0)--(0,2.1)--(-2.1,0)--cycle;
    \fill[ExampleGold!10] (0,0)--(-2.1,0)--(0,-2.1)--cycle;
    \fill[PredictionTeal!9] (0,0)--(0,-2.1)--(2.1,0)--cycle;
    \draw[very thick] (-2.4,0)--(2.4,0);
    \draw[very thick] (0,-2.4)--(0,2.4);
    \node[font=\bfseries] at (0,2.95) {\(\PP^1\times\PP^1\)};
    \node[below] at (2.15,0) {\(e_1\)};
    \node[left] at (0,2.15) {\(e_2\)};
    \node[below] at (-2.15,0) {\(-e_1\)};
    \node[left] at (0,-2.15) {\(-e_2\)};
  \end{scope}

  % weighted projective plane
  \begin{scope}[xshift=-5.4cm,yshift=-3.3cm]
    \fill[WarningRed!8] (0,0)--(2.2,0)--(-1.45,2.15)--cycle;
    \fill[SurfaceGreen!8] (0,0)--(-1.45,2.15)--(0,-2.2)--cycle;
    \fill[DictionaryBlue!8] (0,0)--(0,-2.2)--(2.2,0)--cycle;
    \draw[very thick,DictionaryBlue,-{Latex[length=2.2mm]}] (0,0)--(2.5,0);
    \draw[very thick,WarningRed,-{Latex[length=2.2mm]}] (0,0)--(-1.7,2.5);
    \draw[very thick,SurfaceGreen,-{Latex[length=2.2mm]}] (0,0)--(0,-2.5);
    \node[font=\bfseries] at (0.2,3.0) {\(\PP(1,1,n)\)};
    \node[below] at (2.3,0) {\(u_0=(1,0)\)};
    \node[above left] at (-1.55,2.25) {\(u_1=(-1,n)\)};
    \node[left] at (0,-2.3) {\(u_2=(0,-1)\)};
    \node[WarningRed,fill=white,rounded corners,inner sep=2pt]
      at (0.0,1.0) {\(|\det(u_0,u_1)|=n\)};
  \end{scope}

  % Hirzebruch
  \begin{scope}[xshift=4.8cm,yshift=-3.3cm]
    \fill[DictionaryBlue!8] (0,0)--(2.2,0)--(0,2.2)--cycle;
    \fill[SurfaceGreen!8] (0,0)--(0,2.2)--(-1.75,1.75)--cycle;
    \fill[ExampleGold!10] (0,0)--(-1.75,1.75)--(0,-2.2)--cycle;
    \fill[PredictionTeal!9] (0,0)--(0,-2.2)--(2.2,0)--cycle;
    \draw[very thick,-{Latex[length=2.2mm]}] (0,0)--(2.5,0);
    \draw[very thick,-{Latex[length=2.2mm]}] (0,0)--(0,2.5);
    \draw[very thick,-{Latex[length=2.2mm]}] (0,0)--(-2.0,2.0);
    \draw[very thick,-{Latex[length=2.2mm]}] (0,0)--(0,-2.5);
    \node[font=\bfseries] at (0,3.0) {\(\mathbb F_r\)};
    \node[below] at (2.3,0) {\((1,0)\)};
    \node[right] at (0,2.25) {\((0,1)\)};
    \node[above left] at (-1.85,1.85) {\((-1,r)\)};
    \node[right] at (0,-2.3) {\((0,-1)\)};
  \end{scope}
\end{tikzpicture}
\caption{Four standard complete surface fans. The first, second, and
fourth are smooth. The weighted projective fan is simplicial but has
one nonsmooth maximal cone when \(n>1\).}
\label{fig:four-standard-surface-fans}
\end{figure}

The first two fans are already familiar. Their primitive rays are
\[
\begin{array}{c|c}
\PP^2
&
(1,0),\ (0,1),\ (-1,-1)
\\
\PP^1\times\PP^1
&
(1,0),\ (0,1),\ (-1,0),\ (0,-1).
\end{array}
\]
Every adjacent determinant has absolute value \(1\), and both supports
fill \(\RR^2\). Thus both surfaces are smooth and complete.

\noindent
\textit{Reference.}
These are Cox--Little--Schenck, Examples~3.1.9 and~3.1.12.

\subsection{A weighted projective plane}

Fix \(n\geq2\), and consider the three primitive rays generated by
\[
u_0=(1,0),
\qquad
u_1=(-1,n),
\qquad
u_2=(0,-1).
\]
They satisfy the weighted relation
\[
u_0+u_1+n u_2=0.
\]
Taking the cones generated by each proper subset gives the standard
fan of
\[
\PP(1,1,n).
\]

The fan is complete and simplicial. Two maximal cones are smooth:
\[
|\det(u_1,u_2)|=1,
\qquad
|\det(u_2,u_0)|=1.
\]
The third has
\[
|\det(u_0,u_1)|=n,
\]
so the corresponding torus-fixed point is singular when \(n>1\).

This is the global version of the cyclic quotient phenomenon from
Chapter~2: the real cone is simple, but the primitive generators span
a sublattice of index \(n\).

\begin{example}
For \(n=2\), the singular affine chart is the quadratic cone
\[
V(ac-b^2)\subseteq\AAA^3.
\]
Thus \(\PP(1,1,2)\) is complete and has exactly one singular
torus-fixed point.
\end{example}

\noindent
\textit{Reference.}
Weighted projective fans are Cox--Little--Schenck,
Example~3.1.17. The special case \(\PP(1,1,2)\) and its quadratic-cone
chart are worked out there explicitly.

\begin{pauseandthink}
Use the orbit--cone correspondence to count the torus orbits of
\(\PP(1,1,n)\). Which orbit contains the singular point, and which
cone detects it?
\end{pauseandthink}

\subsection{Hirzebruch surfaces}

For \(r\geq0\), let \(\Sigma_r\) be the complete fan whose rays, in
cyclic order, are
\[
v_1=(1,0),
\qquad
v_2=(0,1),
\qquad
v_3=(-1,r),
\qquad
v_4=(0,-1).
\]
The associated toric surface is the \emph{\(r\)-th Hirzebruch
surface}
\[
\mathbb F_r=X_{\Sigma_r}.
\]

Every adjacent determinant has absolute value \(1\):
\[
\det(v_1,v_2)
=
\det(v_2,v_3)
=
\det(v_3,v_4)
=
\det(v_4,v_1)
=
1.
\]
Hence every \(\mathbb F_r\) is smooth. The fan is complete, so every
\(\mathbb F_r\) is complete.

The lattice projection
\[
p:N=\ZZ^2\longrightarrow\ZZ,
\qquad
p(a,b)=a,
\]
is compatible with the fan of \(\PP^1\). It therefore induces a toric
morphism
\[
\mathbb F_r\longrightarrow\PP^1.
\]
This is the ruling of the Hirzebruch surface.

When \(r=0\), the four rays are the coordinate rays, so
\[
\mathbb F_0\simeq\PP^1\times\PP^1.
\]

\noindent
\textit{Reference.}
This family is Cox--Little--Schenck, Example~3.1.16.

\begin{surfacelaboratory}
For the fan of \(\mathbb F_r\):

\begin{enumerate}[label=(\alph*),wide=0pt]
\item Verify the four determinant calculations.

\item Check compatibility of
\[
p(a,b)=a
\]
with the two rays in the fan of \(\PP^1\).

\item Count the torus orbits of dimensions \(2,1,0\).

\item Compare the fans of \(\mathbb F_0\) and
\(\PP^1\times\PP^1\). Why are they literally the same fan?
\end{enumerate}
\end{surfacelaboratory}

\subsection{Blowing up a torus-fixed point}

The blowup of \(\AAA^2\) from Section~3.3 is the local model for a
general operation.

Let
\[
\sigma=\Cone{u_1,u_2}
\]
be a smooth two-dimensional cone in a fan. Insert the ray
\[
\rho=\Cone{u_1+u_2}
\]
and replace \(\sigma\) by
\[
\Cone{u_1,u_1+u_2},
\qquad
\Cone{u_1+u_2,u_2}.
\]
This operation is the \emph{star subdivision} of the fan at
\(\sigma\).

\begin{proposition}[Star subdivision and blowup]
\label{prop:surface-star-subdivision-blowup}
The toric morphism from the subdivided fan to the original fan is the
blowup of the torus-fixed point \(O(\sigma)\).
\end{proposition}

\noindent
\textit{Reference.}
This is the two-dimensional case of Cox--Little--Schenck,
Proposition~3.3.15.

\begin{figure}[ht]
\centering
\begin{tikzpicture}[scale=0.9]
  % P2
  \begin{scope}[xshift=-4.5cm]
    \fill[DictionaryBlue!9] (0,0)--(2.3,0)--(0,2.3)--cycle;
    \fill[SurfaceGreen!9] (0,0)--(0,2.3)--(-1.7,-1.7)--cycle;
    \fill[ExampleGold!10] (0,0)--(-1.7,-1.7)--(2.3,0)--cycle;
    \draw[very thick,-{Latex[length=2.3mm]}] (0,0)--(2.6,0);
    \draw[very thick,-{Latex[length=2.3mm]}] (0,0)--(0,2.6);
    \draw[very thick,-{Latex[length=2.3mm]}] (0,0)--(-2.0,-2.0);
    \node[font=\bfseries] at (0.3,3.15) {\(\PP^2\)};
    \node[WarningRed,fill=white,rounded corners,inner sep=2pt]
      at (0.8,0.8) {blow up this fixed point};
  \end{scope}

  \draw[very thick,-{Latex[length=2.6mm]}] (-0.7,0.5)--(0.7,0.5);
  \node[above] at (0,0.5) {insert \(e_1+e_2\)};

  % blowup
  \begin{scope}[xshift=4.5cm]
    \fill[DictionaryBlue!9] (0,0)--(2.3,0)--(1.15,1.15)--cycle;
    \fill[PredictionTeal!10] (0,0)--(1.15,1.15)--(0,2.3)--cycle;
    \fill[SurfaceGreen!9] (0,0)--(0,2.3)--(-1.7,-1.7)--cycle;
    \fill[ExampleGold!10] (0,0)--(-1.7,-1.7)--(2.3,0)--cycle;
    \draw[very thick,-{Latex[length=2.3mm]}] (0,0)--(2.6,0);
    \draw[very thick,PredictionTeal,-{Latex[length=2.3mm]}] (0,0)--(2.0,2.0);
    \draw[very thick,-{Latex[length=2.3mm]}] (0,0)--(0,2.6);
    \draw[very thick,-{Latex[length=2.3mm]}] (0,0)--(-2.0,-2.0);
    \node[PredictionTeal,above right] at (1.8,1.8) {new ray};
    \node[font=\bfseries] at (0.2,3.15)
      {\(\Bl_p(\PP^2)\simeq\mathbb F_1\)};
  \end{scope}
\end{tikzpicture}
\caption{Blowing up a torus-fixed point of \(\PP^2\) inserts one ray
and replaces one maximal cone by two. The resulting fan is equivalent
by a lattice change to the fan of \(\mathbb F_1\).}
\label{fig:p2-blowup-f1}
\end{figure}

For the fixed point corresponding to
\[
\Cone{e_1,e_2},
\]
the new fan has rays
\[
e_1,\qquad
e_1+e_2,\qquad
e_2,\qquad
-e_1-e_2.
\]
The matrix
\[
\begin{pmatrix}
1&-1\\
0&1
\end{pmatrix}
\in\operatorname{GL}(2,\ZZ)
\]
takes these rays to
\[
(1,0),\qquad
(0,1),\qquad
(-1,1),\qquad
(0,-1),
\]
the fan of \(\mathbb F_1\). Hence
\[
\Bl_p(\PP^2)\simeq\mathbb F_1
\]
for a torus-fixed point \(p\).

\begin{prediction}
Each blowup at a torus-fixed point inserts one new ray and one new
torus-invariant boundary curve. What combinatorial patterns might
arise after repeatedly blowing up fixed points of a smooth complete
toric surface?
\end{prediction}

\subsection{A compact recognition table}

\begin{center}
\footnotesize
\setlength{\tabcolsep}{3.5pt}
\renewcommand{\arraystretch}{1.22}
\begin{tabularx}{\linewidth}{@{}
  >{\raggedright\arraybackslash}p{0.15\linewidth}
  >{\raggedright\arraybackslash}X
  >{\centering\arraybackslash}p{0.09\linewidth}
  >{\centering\arraybackslash}p{0.10\linewidth}
  >{\raggedright\arraybackslash}X@{}}
\toprule
\textbf{Surface}
&
\textbf{Rays}
&
\textbf{Smooth}
&
\textbf{Complete}
&
\textbf{Feature to recognize}
\\
\midrule
\(\PP^2\)
&
\(e_1,e_2,-e_1-e_2\)
&
yes
&
yes
&
three rays and three fixed points
\\
\(\PP^1\times\PP^1\)
&
\(\pm e_1,\pm e_2\)
&
yes
&
yes
&
product fan
\\
\(\PP(1,1,n)\)
&
\((1,0),(-1,n),(0,-1)\)
&
only if \(n=1\)
&
yes
&
one cone of determinant \(n\)
\\
\(\mathbb F_r\)
&
\((1,0),(0,1),(-1,r),(0,-1)\)
&
yes
&
yes
&
a toric ruling to \(\PP^1\)
\\
\bottomrule
\end{tabularx}
\end{center}

\begin{guidingquestion}
How do divisor classes, torus-invariant divisors, and line bundles
translate the rays of these fans into intersection theory and
projective embeddings?
\end{guidingquestion}

\clearpage
\section*{What this chapter has accomplished}
\addcontentsline{toc}{section}{What this chapter has accomplished}
\markboth{What this chapter has accomplished}{What this chapter has accomplished}

A single cone gives one affine toric chart. A fan records how many such
charts fit together and how they overlap. The chapter has developed
the global dictionary
\[
\boxed{
\text{fan }\Sigma
\quad\longleftrightarrow\quad
\text{normal separated toric variety }X_\Sigma.
}
\]

More precisely, we have seen that:
\[
\begin{array}{c@{\quad\longleftrightarrow\quad}c}
\sigma\in\Sigma
&
U_\sigma\subseteq X_\Sigma
\\[1mm]
\sigma\cap\tau
&
U_\sigma\cap U_\tau
\\[1mm]
\sigma
&
O(\sigma)
\\[1mm]
\tau\preceq\sigma
&
O(\sigma)\subseteq\overline{O(\tau)}
\\[1mm]
|\Sigma|=N_{\RR}
&
X_\Sigma\text{ complete}
\\[1mm]
\Sigma\text{ smooth}
&
X_\Sigma\text{ smooth}.
\end{array}
\]

The surface laboratory supplied a second thread:
\[
\PP^2,\qquad
\PP^1\times\PP^1,\qquad
\PP(1,1,n),\qquad
\mathbb F_r,\qquad
\Bl_p(\PP^2).
\]
In these examples, rays became boundary curves, maximal cones became
fixed points, and subdivisions became blowups.

The next chapter will attach divisor classes and equations to this
combinatorial skeleton.

\begin{exercises}[Comprehension exercises]
\markboth{Comprehension exercises}{Comprehension exercises}

\item Let
\[
\sigma_1=\Cone{(1,0),(0,1)},
\qquad
\sigma_2=\Cone{(1,1),(-1,2)}.
\]
Compute \(\sigma_1\cap\sigma_2\). Decide whether the two cones and all
of their faces form a fan. If not, add the smallest possible
collection of rays and cones that repairs the problem.

\item Let \(\sigma\subseteq N_{\RR}\) be a strongly convex rational
polyhedral cone, and let
\[
\Sigma_\sigma=\{\tau:\tau\preceq\sigma\}.
\]
Show directly from the gluing construction that
\[
X_{\Sigma_\sigma}\simeq U_\sigma.
\]
In other words, explain why gluing all face charts adds no new points.

\item Let \(\Sigma'\subseteq\Sigma\) be a subfan. Prove that the
canonical map
\[
X_{\Sigma'}\longrightarrow X_\Sigma
\]
identifies \(X_{\Sigma'}\) with the torus-invariant open subset
\[
\bigcup_{\sigma\in\Sigma'}U_\sigma.
\]

\item Consider the fan of \(\PP^2\). Remove one maximal cone and every
face that belongs only to that cone.
\begin{enumerate}[label=(\alph*),wide=0pt]
\item Verify that the remaining collection is a subfan.
\item Identify the corresponding open subvariety of \(\PP^2\).
\item List its torus orbits and their closure relations.
\end{enumerate}

\item Let \(\Sigma'\) be the fan of \(\Bl_0(\AAA^2)\), with rays
generated by
\[
e_1,\qquad e_1+e_2,\qquad e_2.
\]
\begin{enumerate}[label=(\alph*),wide=0pt]
\item Compute the coordinate rings of the two maximal-cone charts.
\item Compute their common open chart.
\item Identify the three one-dimensional orbits and the two
zero-dimensional orbits.
\item Use the orbit--cone correspondence to explain why the closure of
the orbit corresponding to \(e_1+e_2\) is a copy of \(\PP^1\).
\end{enumerate}

\item Let \(\Sigma\) be a fan in a rank-\(3\) lattice with
\[
|\Sigma|=N_{\RR}.
\]
Suppose it has \(6\) rays, \(10\) two-dimensional cones, and \(6\)
maximal three-dimensional cones. Determine the number of torus orbits
of each dimension. Which orbits are closed?

\item For each of the following fans, decide whether it is smooth and
whether it is complete:
\[
\begin{array}{ll}
\text{(a)}&
\Cone{(1,0),(1,2)}\text{ and its faces};
\\[1mm]
\text{(b)}&
\text{the complete fan with rays }
(1,0),(0,1),(-1,0),(0,-1);
\\[1mm]
\text{(c)}&
\text{the complete fan with rays }
(1,0),(-1,3),(0,-1).
\end{array}
\]
For every nonsmooth maximal cone, compute the absolute determinant of
its primitive ray generators.

\item Let \(\Sigma_r\) be the fan of the Hirzebruch surface
\(\mathbb F_r\).
\begin{enumerate}[label=(\alph*),wide=0pt]
\item Verify that the lattice map
\[
p(a,b)=a
\]
is compatible with the fan of \(\PP^1\).
\item Determine the image under \(p\) of every ray and every maximal
cone.
\item Describe the induced toric morphism
\[
\mathbb F_r\longrightarrow\PP^1
\]
on the dense tori.
\end{enumerate}

\item Let \(\Sigma_1,\Sigma_2\) be fans. Prove directly from the
orbit--cone correspondence that
\[
O(\sigma_1\times\sigma_2)
\simeq
O(\sigma_1)\times O(\sigma_2).
\]
Use this to recover the orbit decomposition of
\(\PP^1\times\PP^1\).

\item Let
\[
F:\ZZ^2\longrightarrow\ZZ,
\qquad
F(a,b)=a.
\]
Determine whether \(F\) is compatible with:
\begin{enumerate}[label=(\alph*),wide=0pt]
\item the fan of \(\AAA^2\) and the fan of \(\AAA^1\);
\item the fan of \(\PP^1\times\PP^1\) and the fan of \(\PP^1\);
\item the fan of \(\PP^2\) and the fan of \(\PP^1\).
\end{enumerate}
Whenever it is compatible, describe the induced map on the dense
tori.

\item Start with the fan of \(\PP^2\), and subdivide the cone
\[
\Cone{e_1,e_2}
\]
by the ray through \(e_1+e_2\).
\begin{enumerate}[label=(\alph*),wide=0pt]
\item Check that the resulting fan is smooth and complete.
\item Count its torus orbits.
\item Identify its toric surface in two different ways:
as a blowup of \(\PP^2\), and as a Hirzebruch surface.
\end{enumerate}

\end{exercises}

\clearpage
\section*{Challenge exercise}
\addcontentsline{toc}{section}{Challenge exercise}
\markboth{Challenge exercise}{Challenge exercise}

\noindent
\textbf{Weighted toric blowups and the resolution of an
\(\boldsymbol{A_{n-1}}\) singularity.}

\medskip

Assume throughout that \(\Bbbk=\CC\). Consider the Gorenstein cyclic
quotient
\[
X_n
=
\AAA^2/\mu_n,
\qquad
\zeta\cdot(x,y)=(\zeta x,\zeta^{-1}y).
\]
As in the Chapter~2 challenge,
\[
X_n
\simeq
V(uv-w^n)\subseteq\AAA^3,
\]
the \(A_{n-1}\) surface singularity. Torically, it is the affine
variety of the first-quadrant cone
\[
\sigma=\Cone{e_1,e_2}
\]
in the enlarged lattice
\[
N
=
\ZZ^2+
\ZZ\left(\frac1n(1,n-1)\right).
\]

The goal is to resolve the singularity by inserting the primitive
lattice rays lying between \(e_2\) and \(e_1\). Each insertion is a
toric weighted blowup, and the new rays will produce a chain of
rational exceptional curves.

\begin{warningbox}
In this exercise, a \emph{toric weighted blowup} means the proper
birational toric morphism produced by inserting one primitive
interior ray into a two-dimensional cone.

You may use:
\begin{enumerate}[label=(\roman*),wide=0pt]
\item a refinement with the same support induces a proper birational
toric morphism;
\item the toric variety of the complete one-dimensional fan is
\(\PP^1\);
\item a two-dimensional cone is smooth exactly when its primitive ray
generators form a basis of the lattice.
\end{enumerate}
\end{warningbox}

\begin{figure}[ht]
\centering
\begin{tikzpicture}[scale=0.9]
  % One weighted blowup
  \begin{scope}[xshift=-4.8cm]
    \fill[DictionaryBlue!10] (0,0)--(2.8,0)--(1.55,1.1)--cycle;
    \fill[SurfaceGreen!10] (0,0)--(1.55,1.1)--(0,2.8)--cycle;
    \draw[very thick,-{Latex[length=2.3mm]}] (0,0)--(3.0,0);
    \draw[very thick,-{Latex[length=2.3mm]}] (0,0)--(0,3.0);
    \draw[very thick,PredictionTeal,-{Latex[length=2.3mm]}]
      (0,0)--(2.45,1.74);
    \node[below] at (2.8,0) {\(v_1\)};
    \node[left] at (0,2.8) {\(v_0\)};
    \node[PredictionTeal,above right] at (2.15,1.55)
      {\(w=av_0+bv_1\)};
    \node[font=\bfseries] at (0.8,3.55)
      {one weighted blowup};
  \end{scope}

  % A_n resolution
  \begin{scope}[xshift=4.6cm]
    \draw[very thick,-{Latex[length=2.2mm]}] (0,0)--(3.05,0);
    \draw[very thick,-{Latex[length=2.2mm]}] (0,0)--(0.15,3.05);

    \draw[very thick,PredictionTeal,-{Latex[length=2.2mm]}]
      (0,0)--(0.65,2.9);
    \draw[very thick,PredictionTeal,-{Latex[length=2.2mm]}]
      (0,0)--(1.15,2.65);
    \draw[very thick,PredictionTeal,-{Latex[length=2.2mm]}]
      (0,0)--(1.65,2.25);
    \draw[very thick,PredictionTeal,-{Latex[length=2.2mm]}]
      (0,0)--(2.15,1.65);
    \draw[very thick,PredictionTeal,-{Latex[length=2.2mm]}]
      (0,0)--(2.65,0.85);

    \node[left] at (0.15,2.85) {\(u_0=e_2\)};
    \node[above left] at (0.65,2.75) {\(u_1\)};
    \node[above] at (1.15,2.5) {\(u_2\)};
    \node at (1.75,1.75) {\(\cdots\)};
    \node[above right] at (2.55,0.8) {\(u_{n-1}\)};
    \node[below] at (2.85,0) {\(u_n=e_1\)};

    \node[font=\bfseries] at (1.45,3.55)
      {smooth refinement of \(X_n\)};

    \node[draw=PredictionTeal,circle,inner sep=2pt] (E1)
      at (0.35,-1.0) {\(E_1\)};
    \node[draw=PredictionTeal,circle,inner sep=2pt] (E2)
      at (1.2,-1.0) {\(E_2\)};
    \node at (2.0,-1.0) {\(\cdots\)};
    \node[draw=PredictionTeal,circle,inner sep=2pt] (En)
      at (2.9,-1.0) {\(E_{n-1}\)};
    \draw[thick] (E1)--(E2);
    \draw[thick] (E2)--(1.65,-1.0);
    \draw[thick] (2.35,-1.0)--(En);
    \node at (1.6,-1.65)
      {a chain of \(n-1\) rational curves};
  \end{scope}
\end{tikzpicture}
\caption{A primitive interior ray gives a toric weighted blowup. For
the \(A_{n-1}\) singularity, inserting all intermediate primitive rays
produces a smooth fan and an exceptional chain.}
\label{fig:challenge-an-weighted-blowups}
\end{figure}

\begin{enumerate}[label=\textbf{Hint \arabic*.},wide=0pt]

\item Let
\[
\tau=\Cone{v_0,v_1}
\]
be a smooth two-dimensional cone, and let
\[
w=av_0+bv_1
\]
be a primitive interior lattice vector, where
\[
a,b>0,
\qquad
\gcd(a,b)=1.
\]
Insert the ray \(\rho_w=\Cone{w}\).

Show that the two new maximal cones are
\[
\Cone{v_0,w},
\qquad
\Cone{w,v_1}.
\]
In the basis \(v_0,v_1\), compute
\[
|\det(v_0,w)|=b,
\qquad
|\det(w,v_1)|=a.
\]
Explain why the ordinary blowup corresponds to \(a=b=1\), while
general weights may leave quotient singularities in the two new
charts.

\item Let
\[
\pi_w:Y_w\longrightarrow U_\tau
\]
be the morphism defined by this subdivision. Use the orbit--cone
correspondence to show that:
\begin{enumerate}[label=(\alph*),wide=0pt]
\item \(\pi_w\) is an isomorphism away from the fixed point
\(O(\tau)\);
\item the fiber over \(O(\tau)\) is
\[
E=V(\rho_w);
\]
\item the star of \(\rho_w\) is the complete one-dimensional fan.
\end{enumerate}
Deduce that
\[
E\simeq\PP^1.
\]

\item Return to \(X_n\). Set
\[
f_1=\frac1n(1,n-1),
\qquad
f_2=e_2.
\]
Show that \(f_1,f_2\) form a basis of \(N\), and that
\[
e_1=nf_1-(n-1)f_2.
\]

For \(0\leq i\leq n\), define
\[
u_i
=
i f_1-(i-1)f_2.
\]
Verify the simpler formula
\[
u_i
=
\left(\frac{i}{n},\frac{n-i}{n}\right).
\]
In particular,
\[
u_0=e_2,
\qquad
u_n=e_1.
\]

\item Show that the vectors
\[
u_1,\ldots,u_{n-1}
\]
are the primitive lattice points on the line segment joining \(e_2\)
to \(e_1\), ordered from left to right.

Now compute, in the basis \(f_1,f_2\),
\[
\det(u_{i-1},u_i)=-1
\qquad
(1\leq i\leq n).
\]
Conclude that every cone
\[
\Cone{u_{i-1},u_i}
\]
is smooth.

\item Let \(\Sigma'\) be the fan whose rays are
\[
\rho_i=\Cone{u_i},
\qquad
0\leq i\leq n,
\]
and whose maximal cones are
\[
\Cone{u_{i-1},u_i},
\qquad
1\leq i\leq n.
\]
Show that
\[
|\Sigma'|=\sigma.
\]
Therefore the identity map of \(N\) induces a proper birational toric
morphism
\[
\pi:Y=X_{\Sigma'}\longrightarrow X_n.
\]
Why is \(Y\) smooth? Why does this make \(\pi\) a resolution of
singularities?

Explain how the refinement can be built by inserting
\[
\rho_1,\rho_2,\ldots,\rho_{n-1}
\]
one at a time. Interpret \(\pi\) as a sequence of toric weighted
blowups of torus-fixed points.

\item For \(1\leq i\leq n-1\), let
\[
E_i=V(\rho_i).
\]
Show that the star fan of \(\rho_i\) has two opposite rays. Deduce
\[
E_i\simeq\PP^1.
\]

Use the orbit--cone correspondence to prove
\[
E_i\cap E_j\neq\varnothing
\quad\Longleftrightarrow\quad
|i-j|=1.
\]
When \(|i-j|=1\), identify the unique fixed point in the
intersection.

Conclude that the exceptional fiber is the chain
\[
E_1-E_2-\cdots-E_{n-1}.
\]

\item Verify the relation
\[
u_{i-1}+u_{i+1}=2u_i
\qquad
(1\leq i\leq n-1).
\]
A formula from the intersection theory of smooth toric surfaces,
proved later in the course, says that this relation implies
\[
E_i^2=-2.
\]
Record the intersection matrix
\[
E_i\cdot E_j
=
\begin{cases}
-2,&i=j,\\
1,&|i-j|=1,\\
0,&\text{otherwise}.
\end{cases}
\]

Thus the resolution of
\[
V(uv-w^n)
\]
has an exceptional fiber consisting of \(n-1\) rational
\((-2)\)-curves arranged in a chain.

\item Work out the case \(n=5\). List
\[
u_0,u_1,\ldots,u_5,
\]
draw the refined fan, and draw the corresponding chain
\[
E_1-E_2-E_3-E_4.
\]
Mark the five torus-fixed points of the resolution on your picture.
\end{enumerate}

\begin{tcolorbox}[
  colback=PredictionTeal!5,
  colframe=PredictionTeal,
  title={Beyond the Gorenstein case},
  fonttitle=\bfseries
]
For a general quotient
\(\AAA^2/\tfrac1n(1,q)\), the same strategy gives a chain of rational
curves, but the rays need not be equally spaced. The
Hirzebruch--Jung expansion
\(\frac nq=[\![b_1,\ldots,b_r]\!]\) determines relations
\(u_{i-1}+u_{i+1}=b_i u_i\), and hence
\(E_i^2=-b_i\). In the Gorenstein case \(q=n-1\), all the \(b_i\) are
\(2\).
\end{tcolorbox}

\noindent
\textit{Reference.}
See Cox--Little--Schenck, Theorem~3.4.11 and
Theorems~10.2.3, 10.2.5, and~10.4.4.


\clearpage
\section*{Notes and references}
\addcontentsline{toc}{section}{Notes and references}
\markboth{Notes and references}{Notes and references}

The construction of \(X_\Sigma\), the orbit--cone correspondence,
toric morphisms, and completeness follow Cox--Little--Schenck,
Chapter~3. The surface examples in the final catalogue are drawn from
Examples~3.1.9, 3.1.12, 3.1.16, and 3.1.17.

The blowup examples illustrate the general principle that refinements
of fans give proper birational toric morphisms. The challenge exercise
looks ahead to Cox--Little--Schenck, Chapter~10, where cyclic quotient
surface singularities, Hirzebruch--Jung continued fractions, and their
minimal toric resolutions are developed systematically.

The next chapter turns from the orbit stratification to divisors.
Every ray
\[
\rho\in\Sigma(1)
\]
will determine a torus-invariant prime divisor
\[
D_\rho=V(\rho),
\]
and linear relations among the primitive ray generators will begin to
control divisor classes and intersection numbers.
